\section{Différentielles dans les anneaux de caractéristique p}
Les résultats du présent paragraphe, de nature plus spéciale et technique que
ceux des §§19, 20 et 22, ne serviront qu'exceptionnellement dans le cours du
chap. IV. Leur rôle principal est ici dans la démonstration de trois théorèmes
du §22 (22.3.3, 22.5.8 et 22.7.3), dont le premier et le dernier interviennent
de façon essentielle dans la théorie «fine» des anneaux locaux noethériens du
chap. IV, §7.
\oldpage[0\textsubscript{IV-1}]{154}
\subsection{Systèmes de $p$-générateurs et $p$-bases}
\label{0.21.1-fr}

\begin{env}[21.1.1]
\label{0.21.1.1-fr}
Étant donné un nombre $p$ qui est, soit $0$, soit un nombre premier, nous
dirons qu'un anneau $A$ est de caractéristique $p$ s'il existe un homomorphisme
d'anneaux $P\to A$, où $P$ est le corps premier de caractéristique $p$~; on
notera que cet homomorphisme est alors unique, le composé
$\mathbf Z\to P\to A$ étant l'unique homomorphisme de $\mathbf Z$ dans $A$.
Si $A\neq0$, il revient au même de dire que $A$ contient un corps de
caractéristique $p$, l'image de $P$ étant nécessairement un corps isomorphe à
$P$ (et d'ailleurs le seul sous-corps de $A$ isomorphe à $P$).
\end{env}

\begin{env}[21.1.2]
\label{0.21.1.2-fr}
Si $p>0$, dire que $A$ est de caractéristique $p$ équivaut à dire que, dans
$A$, on a $p\cdot1=0$, ou encore $pA=0$. Si $p=0$, dire que $A$ est de
caractéristique $p$ équivaut à dire que pour tout entier $n\neq0$, $n\cdot1$
est inversible dans $A$. Si $A\neq0$, il ne peut exister qu'un seul $p$
(premier ou $0$) tel que $A$ soit de caractéristique $p$~; cela résulte de
l'identité de Bézout $ap+bq=1$ pour deux nombres premiers distincts $p,q$.
Par contre, l'anneau réduit à $0$ est de caractéristique $p$ pour tout $p$.
\end{env}

\begin{env}[21.1.3]
\label{0.21.1.3-fr}
Si $A$ est de caractéristique $p$, il en est de même de toute algèbre sur
$A$. En particulier, pour tout idéal premier $\mathfrak p$ de $A$, le corps
résiduel de $A$ en $\mathfrak p$ est de caractéristique $p$. Réciproquement,
si $p=0$ et si, pour tout idéal maximal $\mathfrak m$ de $A$, le corps
résiduel de $A$ en $\mathfrak m$ est de caractéristique $0$, il en est de même
de $A$, car pour tout entier $n\neq0$, $n\cdot1$ est alors inversible dans
tous les $A_{\mathfrak m}$, donc aussi dans $A$. Par contre, si $p>0$, un
anneau local peut avoir son corps résiduel de caractéristique $p$ sans être
lui-même de caractéristique $p$, comme le montre l'exemple de l'anneau
premier $\mathbf Z_{p\mathbf Z}$ (19.8.3) ou de l'anneau local artinien
$\mathbf Z/p^n\mathbf Z$ pour $n\geq2$, qui ne contiennent pas de corps.

Notons enfin que, pour un anneau (même réduit), les corps résiduels en ses
idéaux premiers peuvent avoir des caractéristiques différentes, comme le
montre l'exemple de $\mathbf Z$.
\end{env}

\begin{env}[21.1.4]
\label{0.21.1.4-fr}
Dans toute la suite de ce paragraphe, nous supposons fixé un nombre premier
$p$ et tous les anneaux seront supposés de caractéristique $p$, sauf mention
expresse du contraire. Pour un tel anneau $A$, l'application $x\mapsto x^p$
est un endomorphisme de $A$, que l'on note $F_A$~; si $A$ est réduit, $F_A$
est injectif. On pose $A^p=F_A(A)$ (sous-anneau de $A$ formé des $x^p$,
$x\in A$)~; on peut naturellement considérer $A$ comme une $A^p$-algèbre.

On peut aussi considérer $A$ comme une $A$-algèbre au moyen de
l'homomorphisme $F_A:x\mapsto x^p$ de $A$ dans $A$~; autrement dit, il
s'agit de l'algèbre $A^{(p)}$ pour laquelle le produit $\lambda\cdot x$ de
$x\in A$ par un scalaire $\lambda\in A$ est le produit $\lambda^p x$ dans
l'anneau $A$. Il est clair que, pour tout homomorphisme d'anneaux $u:A\to B$,
le couple $(u,u)$ est un di-homomorphisme d'algèbres $A^{(p)}\to B^{(p)}$.
Pour tout $A$-module $E$, on pose $E^{(p)}=E\otimes_A A^{(p)}$, où
$A^{(p)}$ est considéré comme un $(A,A)$-bimodule, la structure de
$A$-module à gauche étant celle qu'on vient de définir et la structure de
$A$-module à droite définie par la multiplication dans $A$~; ainsi, pour
$x\in E$, $\alpha,\beta\in A$, on a
\[
  \alpha(x\otimes\beta)=x\otimes(\alpha\beta),\qquad
  (\alpha x)\otimes\beta=x\otimes\alpha^p\beta.
\]
\end{env}

\oldpage[0\textsubscript{IV-1}]{155}
Posant $x^{(p)}=x\otimes1$, on a donc $(\alpha x)^{(p)}=\alpha^p x^{(p)}$.
Lorsque $F_A$ est injectif, on peut identifier $A^{(p)}$ à l'anneau $A$
considéré comme algèbre sur son sous-anneau $A^p$.

\begin{proposition}[21.1.5]
\label{0.21.1.5-fr}
Soient $A$, $B$ deux anneaux, $u:A\to B$ un homomorphisme.
\begin{enumerate}
  \item[(i)] Toute $A$-dérivation de $B$ dans un $B$-module $L$ est nulle
  dans $A[B^p]$ (donc est une $A[B^p]$-dérivation).
  \item[(ii)] Pour toute sous-$A$-algèbre $A'$ de $A[B^p]$, si
  $j:A'\to B$ est l'injection canonique, l'homomorphisme canonique
  \[
  j_{B/A'/A}:\Omega^1_{B/A}\longrightarrow\Omega^1_{B/A'}
  \]
  est bijectif.
  \item[(iii)] Supposons qu'il existe un entier $s\geq0$ tel que
  $u(A)\supset B^{p^s}$. Alors, dans l'anneau $B\otimes_A B$,
  $\mathfrak J_{B/A}$ est un nilidéal.
\end{enumerate}
\end{proposition}

\begin{proof}
(i) Par récurrence à partir de (20.1.1.1), on déduit, pour toute
$A$-dérivation $D:B\to L$, que $D(x^n)=nx^{n-1}D(x)$, d'où en particulier
$D(x^p)=0$.

(ii) Vu (20.4.8), l'assertion (ii) n'est qu'une traduction de (i), cette
dernière s'écrivant $\operatorname{Der}_{A[B^p]}(B,L)=\operatorname{Der}_A(B,L)$
pour tout $B$-module $L$.

(iii) Pour tout $x\in B$, on a
$(x\otimes1-1\otimes x)^{p^s}=x^{p^s}\otimes1-1\otimes x^{p^s}=0$,
puisque $x^{p^s}\in u(A)$. La conclusion résulte de ce que les éléments
$1\otimes x-x\otimes1$ engendrent $\mathfrak J_{B/A}$ (20.4.4).
\end{proof}

Il résulte aussitôt de (21.1.5) que, pour tout couple d'homomorphismes
d'anneaux $A\to B\to C$, on a
\begin{equation}
\tag{21.1.5.1}
\label{0.21.1.5.1-fr}
\Upsilon_{C/B/A}=\Upsilon_{C/B/A'}
\end{equation}
pour toute sous-$A$-algèbre $A'\subset A[B^p]$ de $B$.

D'autre part, (21.1.5) montre en particulier que, pour les modules de
différentielles «absolues»,
\begin{equation}
\tag{21.1.5.2}
\label{0.21.1.5.2-fr}
\Omega^1_A=\Omega^1_{A/A^p}.
\end{equation}

\begin{corollary}[21.1.6]
\label{0.21.1.6-fr}
Supposons que $B$ soit une $A$-algèbre de type fini et qu'il existe un
entier $s\geq0$ tel que $u(A)\supset B^{p^s}$. Si $\Omega^1_{B/A}=0$,
$u$ est surjectif.
\end{corollary}

\begin{proof}
\oldpage[0\textsubscript{IV-1}]{156}
D'après (21.1.5, (iii)), $\mathfrak J_{B/A}$ est un nilidéal~; en outre,
en vertu de (20.4.4) et de l'hypothèse, $\mathfrak J_{B/A}$ est un idéal de
type fini, donc il est nilpotent~; comme la relation $\Omega^1_{B/A}=0$
signifie que $\mathfrak J_{B/A}=\mathfrak J_{B/A}^2$, on en conclut que
$\mathfrak J_{B/A}=0$, ou encore que $\pi:B\otimes_A B\to B$ est bijectif.
Par ailleurs, tout élément de $B$ ayant sa puissance $p^s$-ième dans $u(A)$,
$B$ est entier sur $A$, et étant de type fini, c'est une $A$-algèbre finie.
On est ainsi ramené à prouver le lemme suivant (où l'on ne suppose pas que
les anneaux considérés sont de caractéristique $p$) :
\end{proof}

\begin{lemma}[21.1.6.1]
\label{0.21.1.6.1-fr}
Soient $R$ un anneau, $S$ une $R$-algèbre finie~; si l'homomorphisme
canonique $\pi:S\otimes_R S\to S$ est bijectif, alors l'homomorphisme
structural $u:R\to S$ est surjectif.
\end{lemma}

\begin{proof}
Il suffit de montrer que, pour tout idéal maximal $\mathfrak m$ de $R$, si
l'on pose $T=R-\mathfrak m$, l'homomorphisme
$u_{\mathfrak m}:R_{\mathfrak m}\to T^{-1}S$ est surjectif (Bourbaki,
\emph{Alg. comm.}, chap.~II, §~3, n$^\circ$~3, th.~1)~; or l'hypothèse
entraîne que l'homomorphisme
$(T^{-1}S)\otimes_{R_{\mathfrak m}}(S/\mathfrak mS)\to T^{-1}S$
est bijectif (0$_I$, 1.3.4), et comme $T^{-1}S$ est une
$R_{\mathfrak m}$-algèbre finie, on voit qu'on peut se borner au cas où
$R$ est un anneau local.

Désignant encore par $\mathfrak m$ son idéal maximal, il suffit de prouver que
$u\otimes1:R/\mathfrak m\to S/\mathfrak mS$ est surjectif, en vertu du
lemme de Nakayama ($S$ étant un $R$-module de type fini)~; comme
l'homomorphisme canonique
$(S/\mathfrak mS)\otimes_{R/\mathfrak m}(S/\mathfrak mS)\to S/\mathfrak mS$
est bijectif et que $S/\mathfrak mS$ est une $(R/\mathfrak m)$-algèbre finie,
on est finalement ramené au cas où $R$ est un corps~; mais alors les rangs de
$S\otimes_R S$ et de $S$ sur $R$ ne peuvent être égaux que si $S$ est de rang
$0$ ou $1$, donc si $u:R\to S$ est surjectif.
\end{proof}

\begin{proposition}[21.1.7]
\label{0.21.1.7-fr}
Soient $A$ un anneau, $B$ une $A$-algèbre, $(x_\alpha)_{\alpha\in I}$ une
famille d'éléments de $B$. Considérons les propriétés suivantes~:
\begin{enumerate}
  \item[a)] $B=A[B^p,(x_\alpha)]$, autrement dit la $A[B^p]$-algèbre
  $B$ est engendrée par la famille $(x_\alpha)_{\alpha\in I}$.
  \item[b)] Le $A[B^p]$-module $B$ est engendré par les monômes
  $\prod_\alpha x_\alpha^{n(\alpha)}$, où $(n(\alpha))_{\alpha\in I}$ est
  une famille d'entiers à support fini telle que $0\leq n(\alpha)<p$ pour
  tout $\alpha\in I$.
  \item[c)] Le $B$-module $\Omega^1_{B/A}$ est engendré par les
  $d_{B/A}(x_\alpha)$ $(\alpha\in I)$.
\end{enumerate}
Alors les propriétés a) et b) sont équivalentes et entraînent c)~; si, en
outre, $B$ est une $A[B^p]$-algèbre de type fini, c) est équivalente à a) et b).
\end{proposition}

\begin{proof}
Il est clair que b) entraîne a), et inversement a) entraîne b), car tout
monôme $\prod_\alpha x_\alpha^{m(\alpha)}$, où les $m(\alpha)$ sont des
entiers $\geq0$ (formant une famille à support fini), peut s'écrire
$(\prod_\alpha x_\alpha^{q(\alpha)})^p\prod_\alpha x_\alpha^{r(\alpha)}$,
en divisant chaque $m(\alpha)$ par $p$, ce qui donne
$m(\alpha)=p q(\alpha)+r(\alpha)$ avec $0\leq r(\alpha)<p$.
Le fait que a) entraîne c) résulte de (21.1.5, (ii)) et de (20.4.7).

Supposons inversement c) vérifiée et que $B$ soit une $A[B^p]$-algèbre de
type fini~; soit $B'$ la sous-$A[B^p]$-algèbre de $B$ engendrée par les
$x_\alpha$~; dans la suite exacte (20.5.7.1)
\[
\Omega^1_{B'/A[B^p]}\otimes_{B'}B\longrightarrow
\Omega^1_{B/A[B^p]}\longrightarrow\Omega^1_{B/B'}\longrightarrow0
\]
l'hypothèse entraîne que la flèche de gauche est surjective (compte tenu de
(21.1.5, (ii)))~; on a donc $\Omega^1_{B/B'}=0$, et comme
$B^p\subset B'\subset B$ et que $B$ est une $B'$-algèbre de type fini,
on a nécessairement $B'=B$ par (21.1.6).
\end{proof}

\begin{remark}[21.1.8]
\label{0.21.1.8-fr}
Lorsque $B$ est un corps, nous démontrerons l'équivalence des propriétés a),
b) et c) sans hypothèse de finitude (21.4.5).
\end{remark}

\begin{definition}[21.1.9]
\label{0.21.1.9-fr}
Soient $A$ un anneau (de caractéristique $p$), $B$ une $A$-algèbre. On dit
qu'une famille $(x_\alpha)_{\alpha\in I}$ d'éléments de $B$ est $p$-libre sur
$A$ (resp. un système de $p$-générateurs de $B$ sur $A$, resp. une $p$-base
de $B$ sur $A$) si la famille des monômes
$\prod_\alpha x_\alpha^{n(\alpha)}$ ($0\leq n(\alpha)<p$,
$(n(\alpha))_{\alpha\in I}$ à support fini) est une famille libre (resp. un
système de générateurs, resp. une base) dans le $A[B^p]$-module $B$.
\end{definition}

On a des définitions correspondantes pour un ensemble $M$ d'éléments de $B$,
en considérant la famille définie par l'injection canonique $M\to B$.
Lorsqu'on prend pour $A$ le corps premier $\mathbf F_p$ (auquel cas
$A[B^p]=B^p$), on omet la mention de $A$ dans la définition précédente (ou
l'on dit encore qu'une famille est «absolument» $p$-libre, resp. un système
«absolu» de $p$-générateurs, resp. une $p$-base «absolue»).

\oldpage[0\textsubscript{IV-1}]{157}
\begin{lemma}[21.1.10]
\label{0.21.1.10-fr}
Soient $C$ une sous-$A$-algèbre de $B$ telle que $B^p\subset C$, $M$ une
partie de $C$, $N$ une partie de $B$.
\begin{enumerate}
  \item[(i)] Si $M$ est un système de $p$-générateurs de $C$ sur $A$ et $N$
  un système de $p$-générateurs de $B$ sur $C$, alors $M\cup N$ est un système
  de $p$-générateurs de $B$ sur $A$.
  \item[(ii)] Supposons que $M$ soit une $p$-base de $C$ sur $A$~; alors,
  pour que $N$ soit $p$-libre sur $C$, il faut et il suffit que $M\cup N$
  soit $p$-libre sur $A$.
\end{enumerate}
\end{lemma}

\begin{proof}
On peut se borner au cas où $B^p\subset A\subset C\subset B$. Alors (i)
est un cas particulier du fait que si $P$ (resp. $Q$) est un système de
générateurs du $A$-module $C$ (resp. du $C$-module $B$), l'ensemble des
produits $xy$, où $x\in P$ et $y\in Q$, est un système de générateurs du
$A$-module $B$.

Conservant les mêmes notations, si $P$ est une base du $A$-module $C$,
dire que $Q$ est une famille libre sur $C$ signifie que la relation
$\sum_{\lambda,\mu}a_{\lambda\mu}x_\lambda y_\mu=0$, où
$a_{\lambda\mu}\in A$, $x_\lambda\in P$, $y_\mu\in Q$ (les $x_\lambda$
(resp. les $y_\mu$) étant deux à deux distincts) équivaut à
$a_{\lambda\mu}=0$ pour tout couple $(\lambda,\mu)$, ou encore à
$\sum_\lambda a_{\lambda\mu}x_\lambda=0$ pour tout $\mu$~; d'où
l'assertion (ii).
\end{proof}

\subsection{$p$-bases et lissité formelle}
\label{0.21.2-fr}

\begin{theorem}[21.2.1]
\label{0.21.2.1-fr}
Soient $B$ un anneau, $A$ un sous-anneau de $B$ tel que $B^p\subset A\subset B$.
Soient $E$ une $A$-algèbre, $u:A\to E$ l'homomorphisme structural,
$q:E\to B$ un $A$-homomorphisme surjectif, $\mathfrak R$ son noyau, et
supposons que $\mathfrak R^p=0$. Alors~:
\begin{enumerate}
  \item[(i)] Pour qu'il existe un $A$-homomorphisme $v:B\to E$ inverse à
  droite de l'homomorphisme $q:E\to B$, il faut et il suffit que
  $u(q(z)^p)=z^p$ pour tout $z\in E$.
  \item[(ii)] Lorsque la condition de (i) est satisfaite, pour toute famille
  $(z_\alpha)_{\alpha\in I}$ d'éléments de $E$ telle que
  $q(z_\alpha)=x_\alpha$ pour tout $\alpha\in I$, il existe un
  $A$-homomorphisme $v:B\to E$ et un seul tel que $v(x_\alpha)=z_\alpha$
  pour tout $\alpha\in I$, et $v$ est inverse à droite de $q$.
\end{enumerate}
\end{theorem}

\begin{proof}
\oldpage[0\textsubscript{IV-1}]{158}
Considérons d'abord le cas où $E$ est discret, donc $\mathfrak R$ est
nilpotent~; le raisonnement de (19.4.3) permet en outre de supposer
$\mathfrak R^2=0$ et par suite $\mathfrak R^p=0$. Enfin, en considérant
l'image réciproque de l'extension $E$ de $C$ par $\mathfrak R$, on peut se
borner au cas où $C=B$ et où $v$ est l'identité (19.4.4).

Comme l'homomorphisme d'anneaux $F_B:z\mapsto z^p$ s'annule dans
$\mathfrak R$, il se factorise en $E\xrightarrow q B\xrightarrow{F'}E^{(p)}$,
et $F'$, considéré comme homomorphisme de $B$ dans $E^{(p)}$, est un
$A$-homomorphisme par définition de la structure de $A$-algèbre de $E^{(p)}$
à l'aide de l'homomorphisme composé $A\xrightarrow rE\xrightarrow{F_E}E^{(p)}$,
où $r:A\to E$ est l'homomorphisme structural de l'algèbre $E$~; d'ailleurs,
$r:A^{(p)}\to E^{(p)}$ est aussi un $A$-homomorphisme.
Il y a donc un unique $A$-homomorphisme
$f:A^{(p)}\otimes_A B\to E^{(p)}$ tel que les composés de $f$ et des
homomorphismes canoniques $A^{(p)}\to A^{(p)}\otimes_A B$ et
$B\to A^{(p)}\otimes_A B$ soient respectivement $r$ et $F'$.

Or, par hypothèse, on peut identifier $A^{(p)}\otimes_A B$ à $A[B^p]$,
les homomorphismes canoniques s'identifiant respectivement à $u$ et à $F_B$.
On peut donc considérer $E$ comme une $A[B^p]$-algèbre au moyen de $f$~;
par construction, on a $f(q(z)^p)=z^p$ pour tout $z\in E$~; on est donc
dans les conditions d'application de (21.2.1), d'où le théorème dans ce cas.

Pour passer au cas général, considérons un système fondamental
$(\mathfrak J_\lambda)$ d'idéaux ouverts de $E$, et posons
$E_\lambda=E/\mathfrak J_\lambda$, $\mathfrak R_\lambda=(\mathfrak R+\mathfrak J_\lambda)/\mathfrak J_\lambda$,
$C_\lambda=E_\lambda/\mathfrak R_\lambda=E/(\mathfrak R+\mathfrak J_\lambda)$~;
comme $E=\varprojlim E_\lambda$ (0$_I$, 7.2.2). Pour tout couple
$(\alpha,\lambda)$, désignons par $z_{\alpha\lambda}$ l'image canonique de
$z_\alpha$ dans $E_\lambda$, et par $p_\lambda:C\to C_\lambda$ et
$q_\lambda:E_\lambda\to C_\lambda$ les homomorphismes canoniques. Comme
par hypothèse $\mathfrak R_\lambda$ est nilpotent dans $E_\lambda$, la
première partie de la démonstration prouve l'existence et l'unicité d'un
$A$-homomorphisme $w_\lambda:B\to E_\lambda$ tel que
$p_\lambda\circ v=q_\lambda\circ w_\lambda$ et
$w_\lambda(x_\alpha)=z_{\alpha\lambda}$ pour tout $\alpha$. L'unicité des
$w_\lambda$ montre alors aussitôt que $(w_\lambda)$ est un système projectif
d'homomorphismes, et $w=\varprojlim w_\lambda$ répond à la question.
\end{proof}

\begin{corollary}[21.2.2]
\label{0.21.2.2-fr}
Soient $B$ un anneau, $A$ un sous-anneau de $B$ tel que $B^p\subset A\subset B$.
S'il existe une $p$-base de $B$ sur $A$, $B$ est une $A$-algèbre formellement
lisse relativement à $B^p$ (pour les topologies discrètes).
\end{corollary}

\begin{proof}
Soient $E$ une $A$-extension de $B$ par un $B$-module $L$, $q:E\to B$
l'augmentation, $u:A\to E$ l'homomorphisme structural. Dire que $E$ est
$B^p$-triviale signifie qu'il existe un homomorphisme d'anneaux $v:B\to E$
tel que $q(v(b))=b$ et $v(b^p)=u(b^p)$ pour tout $b\in B$. On en déduit,
pour $z\in E$, que $u(q(z)^p)=v(q(z)^p)=(v(q(z)))^p$~; mais, en vertu de
la relation $L^2=0$, on a aussi $L^p=0$, et comme $v(q(z))-z\in L$, on a
$(v(q(z)))^p=z^p$. La condition de (21.2.1), (i), est donc satisfaite et
$E$ est $B^p$-triviale.
\end{proof}

\begin{corollary}[21.2.3]
\label{0.21.2.3-fr}
Soient $B$ un anneau, $A$ un sous-anneau de $B$ tel que $B^p\subset A\subset B$,
et $\left(x_\alpha\right)_{\alpha\in I}$ une $p$-base de $B$ sur $A$. Soit
$L$ un $B$-module. Alors~:
\begin{enumerate}
  \item[(i)] Pour qu'une dérivation $D$ de $A$ dans $L$ se prolonge en une
  dérivation de $B$ dans $L$, il faut et il suffit que $D$ s'annule dans $B^p$.
  \item[(ii)] Si $D$ s'annule dans $B^p$, alors, pour toute famille
  $(y_\alpha)_{\alpha\in I}$ d'éléments de $L$, il existe une dérivation
  $D'$ de $B$ dans $L$ et une seule prolongeant $D$ et telle que
  $D'(x_\alpha)=y_\alpha$ pour tout $\alpha\in I$.
\end{enumerate}
\end{corollary}

\begin{proof}
Étant donnée une dérivation $D$ de $A$ dans $L$, considérons l'anneau
$E=D_B(L)$ et l'homomorphisme $u:A\to E$ défini par $u(a)=(a,D(a))$~; un
$A$-homomorphisme $v:B\to E$ inverse à droite de l'homomorphisme canonique
$q:E\to B$ est alors de la forme $x\mapsto(x,D'(x))$, où $D'$ est une
dérivation de $B$ dans $L$ prolongeant $D$ (20.1.5)~; comme
$u(q(z)^p)=(q(z)^p,D(q(z)^p))$ pour $z\in E$, le corollaire résulte aussitôt
de (21.2.1) appliqué à $E$.
\end{proof}

Il résulte de (21.2.3) que la suite
\begin{equation}
\tag{21.2.3.1}
\label{0.21.2.3.1-fr}
0\to\operatorname{Der}_A(B,L)\to\operatorname{Der}_{B^p}(B,L)
\to\operatorname{Der}_{B^p}(A,L)\to0
\end{equation}
(cf. (20.2.2.1)) est exacte, et que l'application
\begin{equation}
\tag{21.2.3.2}
\label{0.21.2.3.2-fr}
D'\longmapsto(D'(x_\alpha))_{\alpha\in I}
\end{equation}
est un isomorphisme du $B$-module $\operatorname{Der}_A(B,L)$ sur le
$B$-module produit $L^I$ (en faisant $D=0$ dans (21.2.3)).

\begin{corollary}[21.2.4]
\label{0.21.2.4-fr}
Sous les hypothèses de (21.2.3), la suite
\begin{equation}
\tag{21.2.4.1}
\label{0.21.2.4.1-fr}
0\to\Omega^1_{A/B^p}\otimes_A B\to\Omega^1_{B/B^p}
\to\Omega^1_{B/A}\to0
\end{equation}
est exacte et scindée, et la famille $(d_{B/A}(x_\alpha))_{\alpha\in I}$ est
une base du $B$-module $\Omega^1_{B/A}$.
\end{corollary}

\begin{proof}
Cela résulte aussitôt de (21.2.3) et de la formule
$\operatorname{Der}_A(B,L)=\operatorname{Hom}_B(\Omega^1_{B/A},L)$ (20.4.8).
\end{proof}

\begin{corollary}[21.2.5]
\label{0.21.2.5-fr}
Soient $A$ un anneau, $B$ une $A$-algèbre, $(x_\alpha)_{\alpha\in I}$ une
$p$-base de $B$ sur $A$. Alors $(d_{B/A}(x_\alpha))_{\alpha\in I}$ est une
base du $B$-module $\Omega^1_{B/A}$.
\end{corollary}

\begin{proof}
En utilisant (21.1.5), (ii), on se ramène au cas où $A=A[B^p]$ et il suffit
alors d'appliquer (21.2.4).
\end{proof}

\oldpage[0\textsubscript{IV-1}]{159}
\begin{env}[21.2.6]
\label{0.21.2.6-fr}
Soient $A$, $B$ deux anneaux, $u:A\to B$ un homomorphisme d'anneaux~; on a
(avec les notations de (21.1.4)) le diagramme commutatif
\begin{equation}
\tag{21.2.6.1}
\label{0.21.2.6.1-fr}
\xymatrix{
 A^{(p)} \ar[r]^{u} & B^{(p)}\\
 A \ar[u]^{F_A} \ar[r]^{u} & B \ar[u]_{F_B}
}
\end{equation}
On en déduit canoniquement un homomorphisme de $B$-algèbres
\begin{equation}
\tag{21.2.6.2}
\label{0.21.2.6.2-fr}
u\otimes1:A^{(p)}\otimes_A B\longrightarrow B^{(p)}
\end{equation}
dont l'image est l'anneau $A[B^p]$ (qui est une $B$-algèbre pour
l'homomorphisme $F_B:B\to B^p$).
\end{env}

\begin{theorem}[21.2.7]
\label{0.21.2.7-fr}
Soient $A$ un anneau, $B$ une $A$-algèbre, $u:A\to B$ l'homomorphisme
structural. On suppose vérifiées les conditions suivantes~:
\begin{enumerate}
  \item[(i)] L'homomorphisme (21.2.6.2) déduit de $u$ est injectif.
  \item[(ii)] $B$ admet une $p$-base $(x_\alpha)_{\alpha\in I}$ par rapport à $A$.
\end{enumerate}
Alors $B$ est une $A$-algèbre formellement lisse (pour les topologies
discrètes). De façon plus précise, munissons $A$ et $B$ des topologies
discrètes~; soient $E$ une $A$-algèbre topologique admissible
(0$_I$, 7.1.2), $\mathfrak R$ un idéal de définition de $E$, $C=E/\mathfrak R$,
$v:B\to C$ un $A$-homomorphisme, $q:E\to C$ l'augmentation. Alors, pour
toute famille $(z_\alpha)_{\alpha\in I}$ d'éléments de $E$ telle que
$q(z_\alpha)=v(x_\alpha)$ pour tout $\alpha\in I$, il existe un
$A$-homomorphisme $w:B\to E$ et un seul tel que $v=q\circ w$ et
$w(x_\alpha)=z_\alpha$ pour tout $\alpha\in I$.
\end{theorem}

\begin{proof}
\oldpage[0\textsubscript{IV-1}]{160}
Considérons d'abord le cas où $E$ est discret, donc $\mathfrak R$ nilpotent.
On peut se borner au cas où $v$ est surjectif~; en outre, le raisonnement de
(19.4.3) permet de supposer $\mathfrak R^2=0$ et par suite $\mathfrak R^p=0$.
Enfin, en considérant l'image réciproque par $v$ de l'extension $E$ de $C$
par $\mathfrak R$, on peut se borner au cas où $C=B$ et où $v$ est l'identité
(19.4.4). Comme l'homomorphisme d'anneaux $F_B:B\to B^{(p)}$ s'annule
dans $\mathfrak R$, il se factorise en $E\to B\to B^{(p)}$.
L'homomorphisme $F'$ de $B$ dans $E^{(p)}$ est un $A$-homomorphisme par
définition de la structure de $A$-algèbre de $E^{(p)}$ à l'aide du composé
$A\to E\to E^{(p)}$, et l'homomorphisme $A^{(p)}\to E^{(p)}$ est également
un $A$-homomorphisme. Il existe donc un unique $A$-homomorphisme
$f:A^{(p)}\otimes_A B\to E^{(p)}$ dont les restrictions aux deux facteurs
canoniques sont respectivement l'homomorphisme structural et $F'$.
Par (21.2.6.2), on identifie $A^{(p)}\otimes_A B$ à $A[B^p]$~; on peut
donc considérer $E$ comme une $A[B^p]$-algèbre au moyen de $f$, et par
construction $f(q(z)^p)=z^p$ pour tout $z\in E$. La condition (i) de
(21.2.1) est satisfaite, d'où le théorème dans ce cas.

Pour passer au cas général, considérons un système fondamental
$(\mathfrak J_\lambda)$ d'idéaux ouverts de $E$, et posons
$E_\lambda=E/\mathfrak J_\lambda$, $\mathfrak R_\lambda=(\mathfrak R+\mathfrak J_\lambda)/\mathfrak J_\lambda$,
$C_\lambda=E_\lambda/\mathfrak R_\lambda=E/(\mathfrak R+\mathfrak J_\lambda)$~;
comme $E=\varprojlim E_\lambda$ (0$_I$, 7.2.2). Pour tout couple
$(\alpha,\lambda)$, désignons par $z_{\alpha\lambda}$ l'image canonique de
$z_\alpha$ dans $E_\lambda$, par $p_\lambda:C\to C_\lambda$ l'homomorphisme
canonique et par $q_\lambda:E_\lambda\to C_\lambda$ l'homomorphisme canonique.
Comme par hypothèse $\mathfrak R_\lambda$ est nilpotent dans $E_\lambda$, la
première partie de la démonstration prouve l'existence et l'unicité d'un
$A$-homomorphisme $w_\lambda:B\to E_\lambda$ tel que
$p_\lambda\circ v=q_\lambda\circ w_\lambda$ et
$w_\lambda(x_\alpha)=z_{\alpha\lambda}$ pour tout $\alpha$. L'unicité des
$w_\lambda$ montre alors aussitôt que $(w_\lambda)$ est un système projectif
d'homomorphismes, et $w=\varprojlim w_\lambda$ répond à la question.
\end{proof}

\begin{remark}[21.2.8]
\label{0.21.2.8-fr}
La vérification de l'existence et de l'unicité de l'homomorphisme $w:B\to E$
tel que $w(x_\alpha)=z_\alpha$ pour tout $\alpha$ peut se déduire directement
du fait que $B$ est une $A$-algèbre formellement lisse et que les
$d_{B/A}(x_\alpha)$ forment une base du $B$-module $\Omega^1_{B/A}$ (21.2.4),
sans faire intervenir le fait qu'il s'agit d'une $p$-base (de sorte que le
résultat est valable sans supposer les anneaux de caractéristique $p$).
En effet, on peut se borner au cas où $\mathfrak R^2=0$~; comme
$\operatorname{Der}_A(B,\mathfrak R)=\operatorname{Hom}_B(\Omega^1_{B/A},\mathfrak R)$
(20.4.8), il existe une $A$-dérivation $D$, et une seule, de $B$ dans
$\mathfrak R$ telle que $D(x_\alpha)=z'_\alpha$ pour toute famille
$(z'_\alpha)$ d'éléments de $\mathfrak R$~; la conclusion résulte de (20.1.1).
\end{remark}

\subsection{$p$-bases et modules d'imperfection}
\label{0.21.3-fr}

\begin{env}[21.3.1]
\label{0.21.3.1-fr}
Soient $A$, $B$ deux anneaux (de caractéristique $p$), $i:A\to B$,
$j:B\to A$ deux homomorphismes d'anneaux tels que
\begin{align}
\tag{21.3.1.1}\label{0.21.3.1.1-fr}
j(i(a))&=a^p\quad\text{pour tout }a\in A,\\
\tag{21.3.1.2}\label{0.21.3.1.2-fr}
i(j(b))&=b^p\quad\text{pour tout }b\in B.
\end{align}
Le plus souvent, $i$ sera injectif, de sorte que $A$ s'identifiera par
$i:A\to B$ à un sous-anneau de $B$~; une fois cette identification faite,
l'existence de $j$ implique que $B^p\subset A$, et $j$ s'identifie alors à $F_B$.
\end{env}

\begin{env}[21.3.2]
\label{0.21.3.2-fr}
Si $h:A^p\to A$ est l'injection canonique, on a, d'après (21.3.1.1) et
(21.3.1.2), un diagramme commutatif
\begin{equation}
\tag{21.3.2.1}\label{0.21.3.2.1-fr}
\xymatrix{
 B \ar[r]^{j} & A\\
 A \ar[u]^{i} \ar[r]^{F_A} & A^p \ar[u]_{h}
}
\end{equation}
si bien que le couple $(j,F_A)$ peut être considéré comme un di-homomorphisme
de la $A$-algèbre $B$ (pour $i$) dans la $A^p$-algèbre $A$ (pour $h$).
On en déduit un homomorphisme canonique de $A$-modules
\begin{equation}
\tag{21.3.2.2}\label{0.21.3.2.2-fr}
\pi_{B/A}:\Omega^1_{B/A}\otimes_B A_{[j]}\longrightarrow
\Omega^1_{A/A^p}=\Omega^1_A
\end{equation}
(cf. (20.5.4)~; l'identification de $\Omega^1_{A/A^p}$ et du module des
différentielles absolues $\Omega^1_A$ provient de (21.1.5), (ii)).
\begin{align}
\tag{21.3.2.3}\label{0.21.3.2.3-fr}
\Theta_{B/A}&=\Omega^1_{B/A}\otimes_B A_{[j]},\\
\tag{21.3.2.4}\label{0.21.3.2.4-fr}
\Xi_{B/A}&=\operatorname{Ker}(\pi_{B/A}).
\end{align}
\oldpage[0\textsubscript{IV-1}]{161}
De sorte que l'on a la suite exacte
\begin{equation}
\tag{21.3.2.5}\label{0.21.3.2.5-fr}
0\to\Xi_{B/A}\to\Theta_{B/A}\xrightarrow{\pi_{B/A}}\Omega^1_A
\end{equation}
(on remarquera que ces notations peuvent prêter à confusion puisque
$\Theta_{B/A}$ et $\Xi_{B/A}$ dépendent non seulement de $A$ et de $B$,
mais aussi de $i$ et de $j$).
\end{env}

\begin{env}[21.3.3]
\label{0.21.3.3-fr}
Comme, en vertu de (21.3.1.2), on a $F_B=i\circ j$, on peut écrire pour
tout $B$-module $M$ (cf. (21.1.4))
\begin{equation}
\tag{21.3.3.1}\label{0.21.3.3.1-fr}
M^{(p)}=M\otimes_B B^{(p)}=(M\otimes_B A_{[j]})\otimes_A B_{[i]},
\end{equation}
d'où en particulier
\begin{equation}
\tag{21.3.3.2}\label{0.21.3.3.2-fr}
(\Omega^1_{B/A})^{(p)}=\Theta_{B/A}\otimes_A B_{[i]},
\end{equation}
et l'on déduit donc de (21.3.2.2) un homomorphisme canonique de $B$-modules
\begin{equation}
\tag{21.3.3.3}\label{0.21.3.3.3-fr}
\pi_{B/A}\otimes1_B:(\Omega^1_{B/A})^{(p)}\longrightarrow
\Omega^1_A\otimes_A B_{[i]}.
\end{equation}
Compte tenu de (20.5.4.2), l'image de cet homomorphisme est le $B$-module
engendré par les éléments $d_A(j(b))\otimes1$ pour $b\in B$. Leur image
par l'homomorphisme canonique $i_{B/A}:\Omega^1_A\otimes_A B\to\Omega^1_B$
déduit de $i$ (20.5.2) est donc dans le sous-$B$-module de $\Omega^1_B$
engendré par les $d_B(i(j(b)))$ pour $b\in B$~; en vertu de (21.3.1.2),
cette image est nulle. Autrement dit, on a une suite d'homomorphismes
\begin{equation}
\tag{21.3.3.4}\label{0.21.3.3.4-fr}
0\to\Xi_{B/A}\otimes_A B_{[i]}\to(\Omega^1_{B/A})^{(p)}
\to\Upsilon_{B/A}
\end{equation}
qui n'est pas nécessairement exacte, mais où le composé de deux
homomorphismes consécutifs est nul.
\end{env}

\begin{proposition}[21.3.4]
\label{0.21.3.4-fr}
Si $B$ est un $A$-module plat (pour $i$), la suite (21.3.3.4) est exacte.
\end{proposition}
\begin{proof}
Cela résulte de la définition de la platitude. La proposition s'applique en
particulier lorsque $B^p\subset A\subset B$, $i$ et $j$ étant respectivement
l'injection canonique et $F_B$, et que $B$ admet une $p$-base sur $A$ (de sorte
que $B$ est alors un $A$-module libre). Mais même dans ce cas le noyau
$\Xi_{B/A}$ n'est pas nécessairement nul. Toutefois~:
\end{proof}

\begin{proposition}[21.3.5]
\label{0.21.3.5-fr}
Soient $B$ un anneau réduit, $A$ un sous-anneau de $B$ tel que
$B^p\subset A\subset B$. Supposons qu'il existe une $p$-base
$(x_\lambda)_{\lambda\in L}$ de $B$ sur $A$ et une $p$-base
$(y_\mu)_{\mu\in M}$ de $A$ sur $B^p$~; alors l'homomorphisme canonique
(21.3.3.4)
\begin{equation}
\tag{21.3.5.1}\label{0.21.3.5.1-fr}
(\Omega^1_{B/A})^{(p)}\longrightarrow\Upsilon_{B/A}
\end{equation}
est bijectif, et les éléments $d_A(x_\lambda^p)\otimes1$ forment une base
du $B$-module $\Upsilon_{B/A}$.
\end{proposition}

\begin{proof}
\oldpage[0\textsubscript{IV-1}]{162}
Comme $B$ est réduit, $x\mapsto x^p$ est un isomorphisme de structure de
$B$ sur $B^p$~; par transport de structure au moyen de cet isomorphisme,
on voit que $(x_\lambda^p)_{\lambda\in L}$ et les $y_\mu$ forment une $p$-base
de $B^p$ sur $A^p$ (21.1.10)~; on en conclut par (21.1.5), (ii), et (21.2.5)
que les $d_A(x_\lambda^p)$ et les $d_A(y_\mu)$ forment une base du
$A$-module $\Omega^1_A$. Or l'image de $d_A(x_\lambda^p)\otimes1$ par
$i_{B/A}$ est $d_B(x_\lambda^p)=0$~; d'autre part, l'image de $d_A(y_\mu)$
par $i_{B/A}$ est $d_B(i(y_\mu))$, et comme les $y_\mu$ et les $x_\lambda$
forment une $p$-base de $B$ sur $B^p$ (21.1.10), les $d_B(y_\mu)$ sont une
partie d'une base de $\Omega^1_B$ (21.2.5), donc sont linéairement indépendants
sur $B$. On en déduit que le noyau de $i_{B/A}$ a pour base les
$d_A(x_\lambda^p)\otimes1$~; comme ce sont les images par (21.3.5.1) des
éléments correspondants de $(\Omega^1_{B/A})^{(p)}$, et que ces derniers
forment une base du $B$-module $(\Omega^1_{B/A})^{(p)}$ (21.2.5),
l'homomorphisme (21.3.5.1) de $B$-modules est bijectif.
\end{proof}

\begin{remarks}[21.3.6]
\label{0.21.3.6-fr}
\begin{enumerate}
  \item[(i)] Soient $A'$, $B'$ deux anneaux de caractéristique $p$, et
  supposons donnés deux homomorphismes $i':A'\to B'$, $j':B'\to A'$ vérifiant
  (21.3.1.1) et (21.3.1.2), ainsi que des homomorphismes d'anneaux
  $f:A\to A'$, $g:B\to B'$ rendant commutatif le diagramme
  \[
  \xymatrix{
  A'\ar[r]^{i'}&B'\ar[r]^{j'}&A'\\
  A\ar[u]^f\ar[r]^i&B\ar[u]^g\ar[r]^j&A\ar[u]^f
  }
  \]
  Alors l'homomorphisme canonique $\Omega^1_{B/A}\to\Omega^1_{B'/A'}$
  (20.5.4.3) donne un di-homomorphisme canonique
  $\Theta_{B/A}\to\Theta_{B'/A'}$ rendant commutatif le diagramme
  \[
  \xymatrix{
  0\ar[r]&\Xi_{B/A}\ar[r]\ar[d]&\Theta_{B/A}\ar[r]^{\pi_{B/A}}\ar[d]&\Omega^1_A\ar[d]\\
  0\ar[r]&\Xi_{B'/A'}\ar[r]&\Theta_{B'/A'}\ar[r]^{\pi_{B'/A'}}&\Omega^1_{A'}
  }
  \]
  \item[(ii)] Soit $A'$ une $A$-algèbre quelconque et posons
  $B'=B\otimes_A A'$~; alors $i'=i\otimes1:A'\to B'$ et
  $j'=j\otimes1:B'\to A'$ vérifient (21.3.1.1) et (21.3.1.2), et l'on a
  $\Omega^1_{B'/A'}=\Omega^1_{B/A}\otimes_A A'=\Omega^1_{B/A}\otimes_B B'$
  (20.5.5)~; on en déduit un $A'$-isomorphisme canonique
  \begin{equation}
  \tag{21.3.6.1}\label{0.21.3.6.1-fr}
  \Theta_{B/A}\otimes_A A'\xrightarrow{\sim}\Theta_{B'/A'}
  \end{equation}
  et aussi un $B'$-homomorphisme canonique
  \begin{equation}
  \tag{21.3.6.2}\label{0.21.3.6.2-fr}
  (\Omega^1_{B/A})^{(p)}\otimes_B B'\xrightarrow{\sim}
  (\Omega^1_{B'/A'})^{(p)}.
  \end{equation}
\end{enumerate}
\end{remarks}

\subsection{Cas des extensions de corps}
\label{0.21.4-fr}

\begin{env}[21.4.0]
\label{0.21.4.0-fr}
Soient $K$ un corps de caractéristique $p>0$, $k$ un sous-corps de $K$~;
alors l'anneau $k[K^p]$ est égal au corps $k(K^p)$ puisque $k$ est algébrique
sur $K^p$. On pourra donc appliquer les résultats des numéros précédents en
remplaçant partout $A$, $B$ et $A[B^p]$ par $k$, $K$ et $k(K^p)$.
\end{env}

\begin{lemma}[21.4.1]
\label{0.21.4.1-fr}
Soient $k$ un corps de caractéristique $p>0$, $K$ une extension de $k$.
Pour qu'un élément $x\in K$ soit $p$-libre sur $k$, il faut et il suffit que
$x\notin k(K^p)$.
\end{lemma}

\begin{proof}
En effet, $x$ est racine du polynôme $X^p-x^p$ de $k(K^p)[X]$, et l'on sait
(Bourbaki, \emph{Alg.}, chap.~V, §~8, n$^\circ$~1, prop.~1) que si
$x\notin k(K^p)$ ce polynôme est irréductible, de sorte que les éléments
$1,x,\ldots,x^{p-1}$ forment une base du $k(K^p)$-module $k(K^p)(x)$.
\end{proof}

\begin{theorem}[21.4.2]
\label{0.21.4.2-fr}
Soient $k$ un corps de caractéristique $p>0$, $K$ une extension de $k$,
$S$ un système de $p$-générateurs de $K$ sur $k$, $L\subset S$ une partie
$p$-libre sur $k$. Il existe alors une $p$-base $B$ de $K$ sur $k$ telle que
$L\subset B\subset S$. En particulier, toute extension de $k$ admet une
$p$-base sur $k$.
\end{theorem}

\begin{proof}
On peut se borner au cas où $K^p\subset k$. En vertu du théorème de Zorn,
il existe dans $K$ une partie $B$ telle que $L\subset B\subset S$, $p$-libre
sur $k$ et maximale parmi les parties de $S$ ayant ces propriétés. Il suffit
de voir que le sous-corps $K'$ de $K$ engendré par $k$ et $B$ est égal à $K$.
Dans le cas contraire, il existerait $x\in S$ non dans $K'$~; comme
$K^p\subset k$, on a $K'=K'(K^p)$~; donc $x$ serait $p$-libre sur $K'$
(21.4.1), et par suite $B\cup\{x\}$ serait $p$-libre sur $k$ (21.1.10),
contrairement à la définition de $B$. La dernière assertion de (21.4.2)
s'obtient en prenant $L=\varnothing$, $S=K$.
\end{proof}

\begin{corollary}[21.4.3]
\label{0.21.4.3-fr}
Soient $k$ un corps de caractéristique $p>0$, $K$ une extension de $k$.
Pour qu'une famille $(x_\alpha)$ d'éléments de $K$ soit $p$-libre sur $k$,
il faut et il suffit que, pour tout $\alpha$, $x_\alpha$ n'appartienne pas au
corps $K_\alpha$ engendré par $k(K^p)$ et par les $x_\beta$ d'indice
$\beta\neq\alpha$.
\end{corollary}

\begin{proof}
La condition est nécessaire en vertu de (21.1.10). Inversement, supposons-la
remplie~; on peut se borner au cas où $(x_\alpha)$ est un système de
$p$-générateurs de $K$. Il existe alors une sous-famille de $(x_\alpha)$ qui
est une $p$-base de $K$ (21.4.2), mais cette famille ne peut être distincte de
$(x_\alpha)$, sans quoi, en vertu de l'hypothèse, ce ne serait pas une famille
de $p$-générateurs de $K$.
\end{proof}

\oldpage[0\textsubscript{IV-1}]{163}
\begin{corollary}[21.4.4]
\label{0.21.4.4-fr}
Soient $k$ un corps de caractéristique $p>0$, $K$ une extension de $k$.
Pour que $\Omega^1_{K/k}=0$, il faut et il suffit que $K=k(K^p)$.
En particulier, pour que $\Omega^1_K=0$, il faut et il suffit que $K$ soit
un corps parfait.
\end{corollary}

\begin{proof}
En effet, si $(x_\alpha)$ est une $p$-base de $K$ sur $k$, les
$d_{K/k}(x_\alpha)$ forment une base de l'espace vectoriel $\Omega^1_{K/k}$
(21.2.5).
\end{proof}

\begin{theorem}[21.4.5]
\label{0.21.4.5-fr}
Soient $k$ un corps de caractéristique $p>0$, $K$ une extension de $k$,
$(x_\alpha)$ une famille d'éléments de $K$. Pour que $(x_\alpha)$ soit
$p$-libre sur $k$ (resp. une famille de $p$-générateurs de $K$ sur $k$,
resp. une $p$-base de $K$ sur $k$), il faut et il suffit que la famille
$(d_{K/k}(x_\alpha))$ soit une famille libre (resp. un système de
générateurs, resp. une base) dans l'espace vectoriel $\Omega^1_{K/k}$.
\end{theorem}

\begin{proof}
\oldpage[0\textsubscript{IV-1}]{164}
Soit $K'$ le sous-corps (égal au sous-anneau) de $K$ engendré par
$k(K^p)$ et les $x_\alpha$~; compte tenu de (21.4.3), il suffit de voir que,
pour chaque $\alpha$, $x_\alpha$ n'appartient pas à $K_\alpha$. Or, dans la
suite exacte
\[
\Omega^1_{K_\alpha/k(K^p)}\otimes_{K_\alpha}K\longrightarrow
\Omega^1_{K/k(K^p)}\longrightarrow\Omega^1_{K/K_\alpha}\longrightarrow0,
\]
les images par la flèche de gauche des $d_{K_\alpha/k}(x_\beta)$ pour
$\beta\neq\alpha$ sont les $d_{K/k}(x_\beta)$, donc engendrent un sous-espace
vectoriel de $\Omega^1_{K/k(K^p)}$ ne contenant pas $d_{K/k}(x_\alpha)$~; comme
ce sous-espace est le noyau de la flèche de droite, on voit que
$d_{K/K_\alpha}(x_\alpha)\neq0$, donc $x_\alpha\notin K_\alpha$.

Si maintenant $(x_\alpha)$ est une famille $p$-libre, elle est une sous-famille
d'une $p$-base de $K$ sur $k$ (21.4.2)~; les $d_{K/k}(x_\alpha)$ font donc
partie d'une base de l'espace vectoriel $\Omega^1_{K/k}$ (21.2.5), et sont
par suite linéairement indépendants sur $K$. Réciproquement, si les
$d_{K/k}(x_\alpha)$ sont linéairement indépendants sur $K$, les deux
paragraphes précédents montrent que $(x_\alpha)$ est $p$-libre sur $k$.
\end{proof}

\begin{corollary}[21.4.6]
\label{0.21.4.6-fr}
Soient $k$ un corps de caractéristique $p>0$, $K$ une extension de $k$,
$x$ un élément de $K$. Les trois conditions suivantes sont équivalentes~:
\begin{enumerate}
  \item[a)] $x\notin k(K^p)$.
  \item[b)] $d_{K/k}(x)\neq0$.
  \item[c)] L'élément $x$ est $p$-libre sur $k$.
\end{enumerate}
\end{corollary}

\begin{proposition}[21.4.7]
\label{0.21.4.7-fr}
Soient $L$ un corps de caractéristique $p>0$, $K$ un sous-corps de $L$ tel
que $L^p\subset K\subset L$. Alors la suite d'espaces vectoriels sur $L$
\begin{equation}
\tag{21.4.7.1}\label{0.21.4.7.1-fr}
0\to\Omega^1_{K/L^p}\otimes_KL\to\Omega^1_{L/L^p}
\to\Omega^1_{L/K}\to0
\end{equation}
est exacte~; en d'autres termes, on a $\Upsilon_{L/K/L^p}=0$.
\end{proposition}
\begin{proof}
C'est un cas particulier de (21.2.4), compte tenu de (21.4.2).
\end{proof}

\begin{proposition}[21.4.8]
\label{0.21.4.8-fr}
Sous les hypothèses de (21.4.7), l'homomorphisme canonique
\begin{equation}
\tag{21.4.8.1}\label{0.21.4.8.1-fr}
(\Omega^1_{L/K})^{(p)}\longrightarrow\Upsilon_{L/K}
\end{equation}
est bijectif~; si $(x_\lambda)$ est une $p$-base de $L$ sur $K$, les éléments
$d_K(x_\lambda^p)\otimes1$ forment une base de l'espace vectoriel $\Upsilon_{L/K}$.
\end{proposition}
\begin{proof}
C'est un cas particulier de (21.3.5).
\end{proof}

\subsection{Application~: critères de séparabilité}
\label{0.21.5-fr}

Dans ce numéro et les deux suivants, on ne suppose plus que les anneaux
considérés soient de caractéristique $p>0$.

\begin{env}[21.5.1]
\label{0.21.5.1-fr}
Notons d'abord que le critère (21.2.7) permet de démontrer une partie du
théorème de Cohen sur les extensions séparables (19.6.1), savoir que si $k$
est de caractéristique $p>0$ et si $K$ est une extension séparable de $k$,
alors $K$ est une $k$-algèbre formellement lisse. En effet, $K$ admet une
$p$-base sur $k$ (21.4.2), et d'autre part, il résulte du critère de MacLane
(Bourbaki, \emph{Alg.}, chap.~V, §~8, n$^\circ$~2, prop.~3) que, dans une
clôture algébrique de $K$, $k^{p^{-1}}$ et $K$ sont linéairement disjoints sur
$k$, et par suite que l'homomorphisme canonique
$k^{p^{-1}}\otimes_k K\to k^{p^{-1}}(K)$ est bijectif, ce qui est précisément
la condition (i) de (21.2.7), après transport de structure par
l'isomorphisme $x\mapsto x^p$.
\end{env}
\begin{proposition}[21.5.2]
\label{0.21.5.2-fr}
\textnormal{(MacLane).}
Soient $A$, $B$ deux anneaux de valuation discrète complets, $\mathfrak m$,
$\mathfrak n$ leurs idéaux maximaux respectifs, $K=A/\mathfrak m$ et
$L=B/\mathfrak n$ leurs corps résiduels, $u:A\to B$ un homomorphisme tel que
$Bu(\mathfrak m)=\mathfrak n$, $u_0=u\otimes 1:K\to L$ l'homomorphisme
correspondant pour les corps résiduels. On considère les conditions suivantes~:
\begin{enumerate}
  \item[\rm a)] $L$ est une extension séparable de $K$ (pour $u_0$).
  \item[\rm b)] Pour tout anneau de valuation discrète complet $B'$ d'idéal
  maximal $\mathfrak n'$ et de corps résiduel $L'$,
\oldpage[0\textsubscript{IV-1}]{165}
  tout homomorphisme
  $u':A\to B'$ tel que $B'u'(\mathfrak m)=\mathfrak n'$ et tout
  $K$-isomorphisme $\sigma:L\to L'$ (relatif à $u_0$ et
  $u'_0=u'\otimes1:K\to L'$), il existe un isomorphisme $w:B\to B'$ tel que
  $u'=w\circ u$ et que $w_0=w\otimes1:L\to L'$ soit égal à $\sigma$.
  \item[\rm b$'$)] Pour tout homomorphisme $u':A\to B$ tel que
  $Bu'(\mathfrak m)=\mathfrak n$ et tel que $u'_0=u'\otimes1:K\to L$ soit égal
  à $u_0$, il existe un automorphisme $w$ de $B$ tel que $u'=w\circ u$ et que
  $w_0=w\otimes1:L\to L$ soit l'identité.
  \item[\rm c)] Désignant par $u_1:A/\mathfrak m^2\to B/\mathfrak n^2$
  l'homomorphisme déduit de $u$ par passage aux quotients, alors, pour tout
  homomorphisme local $u'_1:A/\mathfrak m^2\to B/\mathfrak n^2$ tel que
  $\operatorname{gr}_0(u'_1)=\operatorname{gr}_0(u_1)(=u_0)$, il existe un
  automorphisme $w_1$ de $B/\mathfrak n^2$ tel que
  $\operatorname{gr}_0(w_1)$ soit l'identité et que $u'_1=w_1\circ u_1$.
  \item[\rm c$'$)] Pour tout homomorphisme local $u'_1:A/\mathfrak m^2\to
  B/\mathfrak n^2$ tel que $\operatorname{gr}_0(u'_1)=\operatorname{gr}_0(u_1)$
  et $\operatorname{gr}_1(u'_1)=\operatorname{gr}_1(u_1)$
  (homomorphisme de $\mathfrak m/\mathfrak m^2$ dans
  $\mathfrak n/\mathfrak n^2$), il existe un automorphisme $w_1$ de
  $B/\mathfrak n^2$ tel que $\operatorname{gr}_0(w_1)$ soit l'identité et que
  $u'_1=w_1\circ u_1$.
\end{enumerate}
 Alors on a les implications $c)\Leftrightarrow c')\Leftrightarrow a)\Rightarrow b)
 \Rightarrow b')$. Si de plus $A$ est un anneau de Cohen, les cinq conditions
précédentes sont équivalentes.

\end{proposition}
\begin{proof}
Les implications $b)\Rightarrow b')$ et $c)\Rightarrow c')$ sont triviales.
Montrons que $a)$ implique $b)$.
L'homomorphisme $u$ fait de $B$ un $A$-module sans torsion, puisqu'il
transforme par hypothèse une uniformisante de $A$ en une uniformisante de $B$~;
il en résulte que $B$ est un $A$-module plat (0$_{\rm I}$, 6.3.4), donc une
$A$-algèbre de Cohen (19.8.1)~; le fait que $a)$ implique $b)$ est alors
 conséquence de (19.8.2, (i)) appliquée à $G=B'$, $\mathfrak J=\mathfrak n'$,
 $B'$ étant considéré comme $A$-algèbre pour $u'$ et l'homomorphisme
$B\to B'/\mathfrak n'$ étant le composé
$B\to B/\mathfrak n\to B'/\mathfrak n'$, qui est un $A$-homomorphisme en
 vertu de l'hypothèse $u'_0=\sigma\circ u_0$. Le même raisonnement prouve que
$a)$ entraîne $c)$, en prenant cette fois $G=B/\mathfrak n^2$,
$\mathfrak J=\mathfrak n/\mathfrak n^2$, $G$ étant considéré comme
$A$-algèbre pour $u'_1$.
Prouvons en second lieu que $c')$ implique $a)$. Les deux homomorphismes
$u_1$ et $u'_1$ sont tels que $u'_1=u_1+D$, où $D$ est une dérivation de
$A/\mathfrak m^2$ dans $\mathfrak n/\mathfrak n^2$ (20.1.1)~; d'ailleurs,
l'hypothèse $\operatorname{gr}_1(u'_1)=\operatorname{gr}_1(u_1)$ signifie
que $D$ est nulle dans $\mathfrak m/\mathfrak m^2$, et peut par suite être
considérée comme une dérivation de $K=A/\mathfrak m$ dans
$\mathfrak n/\mathfrak n^2$, et $\mathfrak n/\mathfrak n^2$ s'identifie
(par choix d'une uniformisante de $B$) au $K$-module $L$ (pour $u_0$).
D'autre part, les conditions imposées à $w_1$ entraînent que
 $\operatorname{gr}_1(w_1)$ est aussi l'identité (puisque
 $u_1(\mathfrak m/\mathfrak m^2)$ engendre
$\mathfrak n/\mathfrak n^2$)~; $w_1$ est donc de la forme
$x\mapsto x+D'(x)$, où $D'$ est cette fois une dérivation de
$B/\mathfrak n^2$ dans le $(B/\mathfrak n^2)$-module
$\mathfrak n/\mathfrak n^2$~; comme $w_1$ est l'identité dans
$\mathfrak n/\mathfrak n^2$, $D'$ peut encore (par l'identification
précédente de $\mathfrak n/\mathfrak n^2$ et de $L$) être considérée comme
une dérivation de $L$ dans $L$~; enfin, la relation $u'_1=w_1\circ u_1$
signifie que $D'$ prolonge $D$. Notons d'autre part que toute dérivation $D$
de $K$ dans $L$ correspond à un homomorphisme $u'_1$ vérifiant les conditions
de $c')$ (20.1.1)~; la condition $c')$ signifie donc que toute dérivation de
$K$ dans $L$ se prolonge en une dérivation de $L$ dans lui-même, c'est-à-dire
(20.6.5) que $L$ est séparable sur $K$.
Prouvons enfin que, lorsque $A$ est un anneau de Cohen, $b')$ entraîne $c')$.
Prouvons d'abord que sous les hypothèses de $c')$, il existe un homomorphisme
$u':A\to B$ qui vérifie les hypothèses de $b')$ et qui, par passage aux quotients,
donne $u'_1:A/\mathfrak m^2\to B/\mathfrak n^2$. En effet, cela résulte de
(19.8.6, (i)) appliqué à l'homomorphisme composé
$v':A\to A/\mathfrak m^2\xrightarrow{u'_1}B/\mathfrak n^2$~;
\oldpage[0\textsubscript{IV-1}]{166}
ce dernier se factorise donc en $A\xrightarrow{u'}B\to B/\mathfrak n^2$ et
les hypothèses sur $\operatorname{gr}_0(u'_1)$ et
$\operatorname{gr}_1(u'_1)$ entraînent
$\operatorname{gr}_0(u')=\operatorname{gr}_0(u)=u_0$ et
$\operatorname{gr}_1(u')=\operatorname{gr}_1(u)$, donc l'image par $u'$ d'une
uniformisante de $A$ est une uniformisante de $B$, et l'on a par suite
$Bu'(\mathfrak m)=\mathfrak n$. Il suffit alors, pour obtenir un automorphisme
$w_1$ répondant à la question, de prendre l'automorphisme déduit par passage
aux quotients de l'automorphisme $w$ de $B$ fourni par l'application de $b')$
à l'homomorphisme $u'$.
\end{proof}
\begin{remarks}[21.5.3]
\label{0.21.5.3-fr}
\begin{enumerate}
  \item[(i)] Les propriétés différentielles des corps permettent de résoudre la
  question de l'unicité du corps de représentants dans un anneau local
  noethérien complet $C$ (19.8.7, (ii)). Soient $\mathfrak J$ l'idéal maximal
  de $C$, $K=C/\mathfrak J$ le corps résiduel de $C$~; on peut se borner au
  cas $\mathfrak J\ne0$. Supposons qu'il existe un homomorphisme $u:K\to C$
  qui, composé avec l'augmentation $C\to K$, donne l'identité~; alors, pour
  que cet homomorphisme soit unique, il faut et il suffit que
  $\Omega^1_K=0$. En effet, la condition $\Omega^1_K=0$ entraîne que $K$ est
  formellement non ramifié sur son corps premier (20.7.4)~; s'il existait un
  second homomorphisme $v\ne u$ répondant à la question, il y aurait un plus
  grand entier $n$ tel que $\mathfrak J^n$ contienne l'ensemble des
  $u(x)-v(x)$ pour $x\in K$~; par passage au quotient, $u$ et $v$ donneraient
  deux homomorphismes distincts $u_1,v_1$ de $K$ dans
  $C/\mathfrak J^{n+1}$, dont les composés avec
  $C/\mathfrak J^{n+1}\to C/\mathfrak J^n$ seraient égaux, ce qui contredit
  la définition (19.10.2). Inversement, supposons que $\Omega^1_K\ne0$~; il
  existe alors une dérivation $D\ne0$ de $K$ dans
  $\mathfrak J/\mathfrak J^2$ (20.4.8), donc un homomorphisme
  $v_1:K\to C/\mathfrak J^2$ tel que, si $u_1:K\to C/\mathfrak J^2$ est
  obtenu par passage au quotient à partir de $u$, on ait $v_1=u_1+D$
  (20.1.1). Si $k$ est le corps premier de $K$, $C$ est une $k$-algèbre et
  $K$ une $k$-algèbre formellement lisse (19.6.1), et les restrictions de
  $u_1$ et $v_1$ à $k$ coïncident, donc $v_1$ se factorise en
  $K\xrightarrow{v}C\to C/\mathfrak J^2$ et l'on a $v\ne u$.
  Rappelons ((20.6.20) et (21.4.4)) que la condition $\Omega^1_K=0$ signifie
  que $K$ est parfait s'il est de caractéristique $\ne0$, et extension
  algébrique de $\mathbf Q$ s'il est de caractéristique $0$.
  \item[(ii)] De la même manière, soient $W$ un anneau de Cohen, $C$ un anneau
  local noethérien complet, soit $\mathfrak J\ne0$ un idéal de $C$ contenu dans
  l'idéal maximal~; pour que la factorisation $W\to C\to C/\mathfrak J$
  dans (19.8.6, (i)) soit unique, il faut et il suffit que le corps résiduel
  $K$ de l'anneau de Cohen $W$ soit tel que $\Omega^1_K=0$. En effet, si
  $\Omega^1_K\ne0$ et si $\mathfrak m$ désigne l'idéal maximal de $W$, il
  suffit de composer une dérivation $D\ne0$ de $K$ dans
  $\mathfrak J/\mathfrak m\mathfrak J$ avec l'augmentation $W\to K$ pour
  obtenir une dérivation $D'\ne0$ de $W$ dans
  $\mathfrak J/\mathfrak m\mathfrak J$ et l'on termine le raisonnement comme
  dans (i) en formant à l'aide de $D'$ un homomorphisme
  $v_1:W\to C/\mathfrak m\mathfrak J$ distinct de l'homomorphisme
  $u_1:W\to C/\mathfrak m\mathfrak J$ déduit de $u$ par passage au quotient~;
  on achève le raisonnement en invoquant cette fois (19.8.6, (i)). Si au
  contraire $\Omega^1_K=0$, l'unicité de $v$ résulte déjà de (i) lorsque
  $W=K$ est un corps de caractéristique $0$. Dans le cas contraire, on a
  $\widehat{\Omega}^1_W=0$~; en effet, on a alors $\mathfrak m=pW$, $p$ étant
  la caractéristique de $K$ (19.8.5), et l'homomorphisme canonique
  $\mathfrak m/\mathfrak m^2\to\Omega^1_W/\mathfrak m\Omega^1_W$ (20.5.11.2)
  est par suite nul. La suite exacte (20.5.12.1) appliquée à $W$ et à
  $K=W/\mathfrak m$ entraîne alors que $\Omega^1_W=\mathfrak m\Omega^1_W$,
  d'où notre assertion. Mais alors (20.7.4) $W$ est formellement non ramifié
  (pour sa topologie adique) sur $\mathbf Z$, et l'unicité de $v$ se prouve
  comme dans (i).
\end{enumerate}
\end{remarks}

\oldpage[0\textsubscript{IV-1}]{167}
\subsection{Corps admissibles pour une extension}
\label{section:0.21.6-fr}

\begin{env}[21.6.1]
\label{0.21.6.1-fr}
Étant donnés quatre corps $k_0\subset k\subset K\subset L$, il résulte de
(20.6.16) et (20.6.17) que l'on a une suite exacte
\begin{equation}
\tag{21.6.1.1}
\label{0.21.6.1.1-fr}
0\to\Upsilon_{K/k/k_0}\otimes_KL\xrightarrow{v}\Upsilon_{L/k/k_0}
\xrightarrow{u}\Upsilon_{L/K/k_0}\xrightarrow{s}\Upsilon_{L/K/k}\to0.
\end{equation}

Lorsqu'on laisse fixes $k_0$, $K$ et $L$ et qu'on fait «varier» le corps
intermédiaire $k$ entre $k_0$ et $K$, on a évidemment
$\Upsilon_{L/K/k_0}=\Upsilon_{L/K/k}$ lorsque $k=k_0$. Lorsque
l'homomorphisme canonique $s$ de (21.6.1.1) est encore bijectif, on dit que
$k$ est un corps $k_0$-admissible pour l'extension $L$ de $K$. L'intérêt de
l'existence, sous certaines conditions, de tels corps $k$, qui soient pourtant
«suffisamment proches» de $K$ (par exemple tels que $[K:k]$ soit fini), est
qu'ils permettent de remplacer dans certaines questions les modules de
différentielles $\Omega^1_{K/k_0}$ et $\Omega^1_{L/k_0}$ (qui peuvent être
«trop grands», par exemple lorsque $k_0$ est le corps premier) par
$\Omega^1_{K/k}$ et $\Omega^1_{L/k}$, plus aisément maniables.

Lorsque $k_0$ est le \emph{corps premier}, on dira «corps admissible» au lieu
de «corps $k_0$-admissible».
 
On posera
\end{env}

\begin{env}[21.6.1.2]
\label{0.21.6.1.2-fr}
\[
\Delta(L/K,k/k_0)=\operatorname{Coker}
(\Upsilon_{K/k/k_0}\otimes_KL\to\Upsilon_{L/k/k_0})
\simeq\operatorname{Ker}(\Upsilon_{L/K/k_0}\to\Upsilon_{L/K/k})
\]
(espace vectoriel sur $L$)~; son rang sera noté $d(L/K,k/k_0)$ et appelé
\emph{défaut de $k_0$-admissibilité de $k$ pour l'extension $L$ de $K$}
(il est évidemment nul si et seulement si $k$ est $k_0$-admissible pour cette
extension). Lorsque $k_0$ est le corps premier, on écrira
$\Delta(L/K,k)$ et $d(L/K,k)$ au lieu de $\Delta(L/K,k/k_0)$ et
$d(L/K,k/k_0)$.
\end{env}

\begin{proposition}[21.6.2]
\label{0.21.6.2-fr}
Soient $k_0\subset k\subset K\subset L$ quatre corps.
\begin{enumerate}
  \item[\rm (i)] Les conditions suivantes sont équivalentes~:
  \begin{enumerate}
    \item[\rm a)] Le corps $k$ est $k_0$-admissible pour l'extension $L$ de $K$
    (autrement dit, l'homomorphisme $\Upsilon_{L/K/k_0}\to
    \Upsilon_{L/K/k}$ est injectif, donc bijectif).
    \item[\rm b)] L'homomorphisme canonique
    $u:\Upsilon_{K/k/k_0}\to\Upsilon_{L/k/k_0}$ est nul.
    \item[\rm c)] L'homomorphisme canonique
    $v:\Upsilon_{K/k/k_0}\otimes_KL\to\Upsilon_{L/k/k_0}$ est surjectif
    (donc bijectif).
    \item[\rm d)] On a $d(L/K,k/k_0)=0$ (ou
    $\Delta(L/K,k/k_0)=0$).
  \end{enumerate}
  \item[\rm (ii)] Les conditions équivalentes de (i) sont vérifiées lorsque
  l'on est dans l'un des cas suivants~: $\alpha)$ $L$ est séparable sur $k$~;
  $\beta)$ $L$ est séparable sur $K$~; $\gamma)$ on a
  $k\subset k_0(K^p)$, en désignant par $p$ l'exposant caractéristique de $k_0$.
\end{enumerate}
\end{proposition}

\begin{proof}
\begin{enumerate}
  \item[(i)] Les assertions résultent trivialement de l'exactitude de la suite
  (21.6.1.1).
  \item[(ii)] Si $L$ est séparable sur $k$, on a
  $\Upsilon_{L/K/k_0}=0$ (20.6.19), donc la condition b) de (i) est remplie~;
  si $L$ est séparable sur $K$, on a
  $\Upsilon_{L/K/k_0}=0$ (20.6.19), donc la condition a) de (i) est remplie~;
  enfin, si l'on a $k\subset k_0(K^p)$, il en résulte que
  $\Upsilon_{L/K/k_0}=\Upsilon_{L/K/k}$ en vertu de (21.1.5.1), et la
  condition a) de (i) est vérifiée.
\end{enumerate}
\end{proof}
\oldpage[0\textsubscript{IV-1}]{168}
\begin{env}[21.6.3]
\label{0.21.6.3-fr}
Supposons que l'on ait un diagramme commutatif de monomorphismes de corps
\[
\xymatrix{
 k'_0\ar[r]&k'\ar[r]&K'\ar[r]&L'\\
 k_0\ar[u]\ar[r]&k\ar[u]\ar[r]&K\ar[u]\ar[r]&L\ar[u]
}
\]
Il résulte alors de (20.6.17, (ii)) que l'on a un homomorphisme canonique
\begin{equation}
\tag{21.6.3.1}
\label{0.21.6.3.1-fr}
\Delta(L/K,k/k_0)\longrightarrow\Delta(L'/K',k'/k'_0)
\end{equation}
avec une propriété de transitivité évidente, de sorte que l'on peut dire que
$\Delta(L/K,k/k_0)$ est un \emph{foncteur} en le quadruplet
$(k_0,k,K,L)$.
\end{env}

\begin{proposition}[21.6.4]
\label{0.21.6.4-fr}
\begin{enumerate}
  \item[\rm (i)] Soient $k\subset k'\subset k''\subset K\subset L$ cinq corps.
  On a une suite exacte d'homomorphismes canoniques
\begin{equation}
\tag{21.6.4.1}
\label{0.21.6.4.1-fr}
0\to\Delta(L/K,k'/k)\to\Delta(L/K,k''/k)\to\Delta(L/K,k''/k')\to0
\end{equation}
et par suite l'égalité
\begin{equation}
\tag{21.6.4.2}
\label{0.21.6.4.2-fr}
d(L/K,k''/k)=d(L/K,k''/k')+d(L/K,k'/k).
\end{equation}
  \item[\rm (ii)] Soient $k_0\subset k\subset K\subset L\subset M$ cinq corps.
  On a une suite exacte d'homomorphismes canoniques
\begin{equation}
\tag{21.6.4.3}
\label{0.21.6.4.3-fr}
0\to\Delta(L/K,k/k_0)\otimes_LM\to\Delta(M/K,k/k_0)
\to\Delta(M/L,k/k_0)\to0
\end{equation}
et par suite l'égalité
\begin{equation}
\tag{21.6.4.4}
\label{0.21.6.4.4-fr}
d(M/K,k/k_0)=d(M/L,k/k_0)+d(L/K,k/k_0).
\end{equation}
\end{enumerate}
\end{proposition}

\begin{proof}
\begin{enumerate}
  \item[(i)] Considérons le diagramme commutatif
\[
\xymatrix{
0\ar[r]&\Upsilon_{K/k'/k}\otimes_KL\ar[r]\ar[d]&
\Upsilon_{K/k''/k}\otimes_KL\ar[r]\ar[d]&
\Upsilon_{K/k''/k'}\otimes_KL\ar[r]\ar[d]&0\\
0\ar[r]&\Upsilon_{L/k'/k}\ar[r]\ar[d]&
\Upsilon_{L/k''/k}\ar[r]\ar[d]&
\Upsilon_{L/k''/k'}\ar[r]\ar[d]&0\\
0\ar[r]&\Delta(L/K,k'/k)\ar[r]&\Delta(L/K,k''/k)\ar[r]&
\Delta(L/K,k''/k')\ar[r]&0
}
\]
où l'on peut considérer les trois lignes comme des \emph{complexes}
$T_1,T_2,T_3$ respectivement~; la suite exacte (21.6.1.1) et la définition
de $\Delta$ (21.6.1.2) montrent que l'on a une suite exacte de complexes
$0\to T_1\to T_2\to T_3\to0$~; appliquons-lui la suite exacte de
\oldpage[0\textsubscript{IV-1}]{169}
cohomologie, et notons que, en vertu de l'exactitude de (21.6.1.1), la
cohomologie de $T_1$ et celle de $T_2$ sont nulles sauf en un seul et même
degré, pour lequel les modules de cohomologie sont tous deux égaux à
$\Upsilon_{k''/k'/k}\otimes_{k''}L$~; comme $T_1$ et $T_2$ ont ainsi même
cohomologie, celle de $T_3$ est nécessairement nulle, ce qui prouve (i).
  \item[(ii)] Considérons de même le diagramme commutatif
\[
\xymatrix{
0\ar[r]&\Delta(L/K,k/k_0)\otimes_LM\ar[r]\ar[d]&
\Delta(M/K,k/k_0)\ar[r]\ar[d]&\Delta(M/L,k/k_0)\ar[r]&0\\
0\ar[r]&\Upsilon_{L/K/k_0}\otimes_LM\ar[r]\ar[d]&
\Upsilon_{M/K/k_0}\ar[r]\ar[d]&\Upsilon_{M/L/k_0}\ar[r]&0\\
0\ar[r]&\Upsilon_{L/K/k}\otimes_LM\ar[r]&
\Upsilon_{M/K/k}\ar[r]&\Upsilon_{M/L/k}\ar[r]&0
}
\]
où de nouveau on considère les trois lignes comme des \emph{complexes}
$T'_1,T'_2,T'_3$~; la suite exacte (21.6.1.1) et la définition de $\Delta$
(21.6.1.2) donnent ici une suite exacte de complexes
$0\to T'_1\to T'_2\to T'_3\to0$ à laquelle on applique encore la suite exacte
de cohomologie~; cette fois, en vertu de l'exactitude de (21.6.1.1), la
cohomologie de $T'_2$ et celle de $T'_3$ sont nulles sauf en un seul et même
degré, pour lequel les modules de cohomologie sont tous deux égaux à
$\Upsilon_{M/L/k_0}$~; on conclut donc ici que la cohomologie de $T'_1$ est
nulle, ce qui établit (ii).
\end{enumerate}
\end{proof}

\begin{corollary}[21.6.5]
\label{0.21.6.5-fr}
\begin{enumerate}
  \item[(i)] Étant donnés cinq corps $k\subset k'\subset k''\subset K\subset L$,
  pour que $k''$ soit $k$-admissible pour l'extension $L$ de $K$, il faut et
  il suffit que $k'$ soit $k$-admissible et que $k''$ soit $k'$-admissible
  pour l'extension $L$ de $K$.
  \item[(ii)] Étant donnés cinq corps
  $k_0\subset k\subset K\subset L\subset M$, pour que $k$ soit
  $k_0$-admissible pour l'extension $M$ de $K$, il faut et il suffit qu'il le
  soit pour l'extension $L$ de $K$ et pour l'extension $M$ de $L$.
\end{enumerate}

Cela résulte aussitôt des relations (21.6.4.2) et (21.6.4.4), les valeurs de
$d$ étant $\geqslant0$.
\end{corollary}

\begin{corollary}[21.6.6]
\label{0.21.6.6-fr}
Soient $k_0\subset k\subset K\subset L$ quatre corps, et supposons que $k$
soit $k_0$-admissible pour l'extension $L$ de $K$. Alors, si $k'_0$, $k'$,
$K'$, $L'$ sont quatre corps tels que
$k_0\subset k'_0\subset k'\subset k\subset K\subset K'\subset L'\subset L$,
$k'$ est $k'_0$-admissible pour l'extension $L'$ de $K'$.
\end{corollary}

\subsection{L'égalité de Cartier}
\label{subsection:0.21.7-fr}

Le résultat suivant traduit en termes de différentielles un théorème de
MacLane sur les dérivations~:

\oldpage[0\textsubscript{IV-1}]{170}
\begin{theorem}[21.7.1]
\label{0.21.7.1-fr}
\textnormal{(Cartier).}
Soient $K$ un corps, $L$ une extension de type fini de $K$. Alors
$\Omega^1_{L/K}$ et $\Upsilon_{L/K}$ sont des $L$-espaces vectoriels de rang
fini, et l'on a
\begin{equation}
\tag{21.7.1.1}
\label{0.21.7.1.1-fr}
\operatorname{rg}\Omega^1_{L/K}-\operatorname{rg}\Upsilon_{L/K}
=\deg.\operatorname{tr}_K L.
\end{equation}
\end{theorem}

\begin{proof}
Si $L'$ est un corps tel que $K\subset L'\subset L$, on a la suite exacte
(20.6.15.1) (appliquée à $\Lambda=P$, corps premier de $K$, $A=K$,
$B=L'$, $C=L$)
\[
0\to\Upsilon_{L'/K}\otimes_{L'}L\to\Upsilon_{L/K}
\to\Omega^1_{L'/K}\otimes_{L'}L\to\Omega^1_{L/K}
\to\Omega^1_{L/L'}\to0
\]
d'où (0$_{\rm III}$, 11.10.2)
\[
\operatorname{rg}_L\Omega^1_{L/K}-\operatorname{rg}_L\Upsilon_{L/K}
=(\operatorname{rg}_{L'}\Omega^1_{L/L'}-\operatorname{rg}_{L'}\Upsilon_{L/L'})
+(\operatorname{rg}_{L'}\Omega^1_{L'/K}-\operatorname{rg}_{L'}\Upsilon_{L'/K}).
\]
Comme par ailleurs $\deg.\operatorname{tr}_K L=\deg.\operatorname{tr}_K L'
+\deg.\operatorname{tr}_{L'}L$, on voit, par récurrence sur le nombre de
générateurs de l'extension $L$, que l'on est ramené à prouver (21.7.1.1)
lorsque $L=K(x)$. Distinguons alors trois cas~:
\begin{enumerate}
  \item[\rm a)] $x$ est transcendant sur $K$~; comme $L$ est séparable sur
  $K$, on a alors $\Upsilon_{L/K}=0$ (20.6.3)~; d'autre part, (20.5.9) et
  (20.4.13) montrent que $\Omega^1_{L/K}$ est de rang $1$, d'où
  (21.7.1.1) dans ce cas.
  \item[\rm b)] $L$ est une extension algébrique séparable de $K$, de sorte
  que l'on a encore $\Upsilon_{L/K}=0$ (20.6.3). D'autre part, on a
  $\Omega^1_{L/K}=0$ par (20.6.20), d'où encore (21.7.1.1).
  \item[\rm c)] $L$ est une extension algébrique inséparable de $K$~; le
  raisonnement du début montre que l'on peut se borner au cas où $L^p\subset K$~;
  il résulte donc de (21.4.8) que l'on a
  $\operatorname{rg}_L\Upsilon_{L/K}=\operatorname{rg}_L\Omega^1_{L/K}$,
  d'où encore (21.7.1.1).
\end{enumerate}
\end{proof}

\begin{corollary}[21.7.2]
\label{0.21.7.2-fr}
Soient $K$ un corps, $L$ une extension de type fini de $K$, $k$ un sous-corps
de $K$. Alors $\Omega^1_{L/K}$ et $\Upsilon_{L/K/k}$ sont des espaces
vectoriels de rang fini sur $L$, et l'on a
\begin{equation}
\tag{21.7.2.1}
\label{0.21.7.2.1-fr}
\operatorname{rg}\Omega^1_{L/K}-\operatorname{rg}\Upsilon_{L/K/k}
=\deg.\operatorname{tr}_K L+d(L/K,k).
\end{equation}
Par suite, on a l'inégalité
\begin{equation}
\tag{21.7.2.2}
\label{0.21.7.2.2-fr}
\operatorname{rg}\Omega^1_{L/K}-\operatorname{rg}\Upsilon_{L/K/k}
\geqslant\deg.\operatorname{tr}_K L.
\end{equation}
En outre, pour que les deux membres de (21.7.2.2) soient égaux, il faut et
il suffit que $k$ soit un corps admissible pour l'extension $L$ de $K$.
\end{corollary}

\begin{proof}
En effet, comme l'homomorphisme
$\Upsilon_{L/K}\to\Upsilon_{L/K/k}$ est surjectif, on a par définition
(21.6.1.2)
$\operatorname{rg}\Upsilon_{L/K}-\operatorname{rg}\Upsilon_{L/K/k}
=d(L/K,k)$, et le corollaire résulte donc aussitôt de (21.7.1) et de
(21.6.2).
\end{proof}

\begin{corollary}[21.7.3]
\label{0.21.7.3-fr}
Soient $k\subset K\subset L$ trois corps tels que $L$ soit une extension de
type fini de $k$. Alors on a
\begin{equation}
\tag{21.7.3.1}
\label{0.21.7.3.1-fr}
\operatorname{rg}_L\Omega^1_{L/k}-\operatorname{rg}_K\Omega^1_{K/k}
=\deg.\operatorname{tr}_K L+d(L/K,k).
\end{equation}
\oldpage[0\textsubscript{IV-1}]{171}
Par suite, on a l'inégalité
\begin{equation}
\tag{21.7.3.2}
\label{0.21.7.3.2-fr}
\operatorname{rg}_L\Omega^1_{L/k}-\operatorname{rg}_K\Omega^1_{K/k}
\geqslant\deg.\operatorname{tr}_K L.
\end{equation}
L'égalité étant atteinte si et seulement si $k$ est un corps admissible pour
l'extension $L$ de $K$.
\end{corollary}

\begin{proof}
On a en effet la suite exacte (20.6.1.1)
\[
0\to\Upsilon_{L/K/k}\to\Omega^1_{K/k}\otimes_KL\to\Omega^1_{L/k}
\to\Omega^1_{L/K}\to0
\]
donc le premier membre de (21.7.3.1) est égal à
\[
\operatorname{rg}_L\Omega^1_{L/k}-\operatorname{rg}_L\Upsilon_{L/K/k}
\]
et il suffit d'appliquer (21.7.2).
\end{proof}

L'intérêt de ce dernier corollaire est qu'il ne fait intervenir que des
modules de différentielles, à l'exclusion de modules d'imperfection. Nous
verrons d'ailleurs plus bas (21.8.6) que pour toute extension de type fini
$L$ de $K$, il existe un sous-corps $k$ de $K$ tel que $[K:k]$ soit fini et
qui est admissible pour l'extension $L$ de $K$, de sorte que les deux membres
de (21.7.3.2) sont alors égaux.

\begin{corollary}[21.7.4]
\label{0.21.7.4-fr}
Soit $K$ une extension de type fini d'un corps $k$. Les conditions suivantes
sont équivalentes~:
\begin{enumerate}
  \item[\rm a)] $K$ est une extension finie séparable de $k$.
  \item[\rm b)] $K$ est une $k$-algèbre formellement non ramifiée (19.10.2).
  \item[\rm c)] $K$ est une $k$-algèbre formellement étale (19.10.2).
  \item[\rm d)] On a $\Omega^1_{K/k}=0$.
\end{enumerate}
\end{corollary}

\begin{proof}
En effet, on sait (les topologies étant discrètes) que les conditions b) et
d) sont équivalentes (20.7.4) et que a) entraîne que $K$ est une
$k$-algèbre formellement lisse (19.6.1), donc la conjonction de a) et b)
équivaut à c). Tout revient à voir que d) entraîne a). Or, en vertu de
(21.7.1.1), la relation $\Omega^1_{K/k}=0$ entraîne que $K$ est algébrique
(donc fini) sur $k$ et que $\Upsilon_{K/k}=0$, c'est-à-dire (20.6.3) que
$K$ est séparable sur $k$.
\end{proof}

\begin{remark}[21.7.5]
\label{0.21.7.5-fr}
En vertu de (20.6.3.2) et de (0$_{\rm III}$, 11.10.2), le premier membre de
(21.7.1.1) n'est autre que la \emph{caractéristique d'Euler-Poincaré} du
complexe $K_\bullet(C/B/A)$ introduit dans (20.6.3), pour $C=L$, $B=K$ et
$A=P$ (corps premier de $K$). Dans le chapitre consacré à la théorie des
intersections et au théorème de Riemann--Roch, un rôle important sera joué
également par les caractéristiques d'Euler-Poincaré généralisées (à valeurs
dans des groupes de classes de $\mathcal O_X$-Modules) de complexes
généralisant les complexes $F_\bullet(C/A)$ considérés dans (20.6.22).
\end{remark}

\subsection{Critères d'admissibilité}
\label{subsection:0.21.8-fr}

Nous revenons à nos conventions antérieures et supposons donc que tous les
corps considérés dans ce numéro sont de caractéristique $p>0$.

\begin{lemma}[21.8.1]
\label{0.21.8.1-fr}
Soient $K$ un corps, $k$ un sous-corps de $K$, $(k_\alpha)_{\alpha\in I}$ une
famille filtrante décroissante de sous-corps de $K$ telle que
$k=\bigcap_{\alpha\in I}k_\alpha$. Soient $V$ un espace vectoriel sur $K$,
$(a_i)_{1\leqslant i\leqslant n}$ une famille finie de vecteurs de $V$~; si la
famille $(a_i)$ est libre sur $k$, il existe un indice $\gamma$ tel qu'elle
soit aussi libre sur $k_\gamma$.
\end{lemma}

\oldpage[0\textsubscript{IV-1}]{172}
\begin{proof}
Soit $r$ le rang de la famille $(a_i)_{1\leqslant i\leqslant n}$ sur $K$, et
raisonnons par récurrence sur $n-r$~; la proposition est évidente pour $n=r$,
car alors la famille $(a_i)_{1\leqslant i\leqslant r}$ est libre sur $K$, donc
sur tout sous-corps de $K$. Supposons par exemple que
$(a_i)_{1\leqslant i\leqslant r}$ soit libre sur $K$, et écrivons
$a_{r+1}=\sum_{i=1}^r\lambda_i a_i$ avec $\lambda_i\in K$~; la famille
$(a_i)_{1\leqslant i\leqslant r+1}$ étant libre sur $k$, les $\lambda_i$ ne
peuvent tous appartenir à $k$~; supposons par exemple que
$\lambda_{i_0}\notin k$. Alors il existe un indice $\beta$ tel que
$\lambda_{i_0}\notin k_\beta$~; on en conclut que la famille $(a_i)_{1\leqslant i\leqslant r+1}$ est libre sur $k_\beta$~; en effet, comme la
famille $(a_i)_{1\leqslant i\leqslant r}$ est libre sur tout sous-corps de
$K$, si la famille $(a_i)_{1\leqslant i\leqslant r+1}$ n'était pas libre sur
$k_\beta$, $a_{r+1}$ serait égal à une combinaison linéaire des $a_i$ à
coefficients dans $k_\beta$, et comme ces coefficients sont nécessairement
les $\lambda_i$, on aboutit à une contradiction. Il suffit maintenant
d'appliquer l'hypothèse de récurrence en remplaçant $K$ par $k_\beta$ et la
famille $(k_\alpha)_{\alpha\in I}$ par la sous-famille des $k_\alpha$
contenus dans $k_\beta$.
\end{proof}

\begin{lemma}[21.8.2]
\label{0.21.8.2-fr}
Soient $K$ un corps, $p$ son exposant caractéristique, $k_0$ un sous-corps de
$K$, $(k_\alpha)_{\alpha\in I}$ une famille filtrante décroissante de
sous-corps de $K$ telle que $\bigcap_{\alpha\in I}k_\alpha(K^p)=k_0(K^p)$.
Si $(x_i)_{1\leqslant i\leqslant n}$ est une famille finie d'éléments de $K$
qui est $p$-libre sur $k_0$, il existe un indice $\alpha$ tel que $(x_i)$ soit
$p$-libre sur $k_\alpha$.

En effet, dire que la famille $(x_i)$ est $p$-libre sur un sous-corps $k$
de $K$ signifie que la famille finie des monômes
$\prod_{i=0}^n x_i^{m(i)}$ avec $0\leqslant m(i)<p$ est libre sur $k(K^p)$~;
il suffit donc d'appliquer le lemme (21.8.1) à cette famille de monômes dans
l'espace vectoriel $V=K$, et aux sous-corps $k_\alpha(K^p)$ et $k_0(K^p)$
de $K$.
\end{lemma}

\begin{theorem}[21.8.3]
\label{0.21.8.3-fr}
Soient $K$ un corps de caractéristique $p>0$, $k_0$ un sous-corps de $K$,
$(k_\alpha)_{\alpha\in I}$ une famille filtrante décroissante de sous-corps
de $K$ contenant $k_0$. Les conditions suivantes sont équivalentes~:
\begin{enumerate}
  \item[\rm a)] $\bigcap_{\alpha\in I}k_\alpha(K^p)=k_0(K^p)$.
  \item[\rm b)] Pour toute extension $L$ de $K$ telle que
  $\Upsilon_{L/K/k_\alpha}$ soit un $L$-espace vectoriel de rang fini
  (ce qui a lieu en particulier si $L$ est une extension de type fini en
  vertu de (21.7.2)), il existe $\alpha\in I$ tel que $k_\alpha$ soit
  $k_0$-admissible pour l'extension $L$ de $K$.
  \item[\rm b$'$)] Pour toute extension $L=K(x)$ de $K$, avec $x^p\in K$,
  il existe $\alpha\in I$ tel que $k_\alpha$ soit $k_0$-admissible pour
  l'extension $L$ de $K$.
  \item[\rm c)] L'application canonique
\begin{equation}
\tag{21.8.3.1}
\label{0.21.8.3.1-fr}
\Omega^1_{K/k_0}\longrightarrow\varprojlim_{\alpha}\Omega^1_{K/k_\alpha}
\end{equation}
est injective.
\end{enumerate}
\end{theorem}

\begin{proof}
L'application canonique (21.8.3.1) est bien entendu obtenue par passage à la
limite projective dans le système projectif d'homomorphismes
$\Omega^1_{K/k_0}\to\Omega^1_{K/k_\alpha}$ (20.5.3.3). Nous allons prouver
le théorème suivant le schéma logique $c)\Rightarrow b)\Rightarrow b')
\Rightarrow a)\Rightarrow c)$.

Dire que $k_\alpha$ est $k_0$-admissible pour l'extension $L$ de $K$ signifie
que l'homomorphisme
\oldpage[0\textsubscript{IV-1}]{173}
canonique $\Upsilon_{L/K/k_0}\to\Upsilon_{L/K/k_\alpha}$ est injectif~; or,
si $N_\alpha$ est le noyau de cet homomorphisme, les $N_\alpha$ forment un
système filtrant décroissant de sous-espaces vectoriels de
$\Upsilon_{L/K/k_0}$, et comme $\Upsilon_{L/K/k_0}$ est par hypothèse de rang
fini, il revient au même de dire que l'un des $N_\alpha$ est nul ou que leur
intersection est nulle. Mais cette intersection n'est autre que le noyau de
l'homomorphisme limite projective
$\Upsilon_{L/K/k_0}\to\varprojlim_\alpha\Upsilon_{L/K/k_\alpha}$. Or, on a
le diagramme commutatif
\[
\xymatrix{
\Upsilon_{L/K/k_0}\ar[r]\ar[d]&\varprojlim\Upsilon_{L/K/k_\alpha}\ar[d]\\
\Omega^1_{K/k_0}\otimes_KL\ar[r]&
\varprojlim(\Omega^1_{K/k_\alpha}\otimes_KL)
}
\]
où la flèche verticale de gauche est injective par définition. Pour prouver
que $c)$ entraîne $b)$, il suffit donc de montrer que l'application canonique
\begin{equation}
\tag{21.8.3.2}
\label{0.21.8.3.2-fr}
\Omega^1_{K/k_0}\otimes_KV\longrightarrow
\varprojlim_{\alpha}(\Omega^1_{K/k_\alpha}\otimes_KV)
\end{equation}
est injective pour tout espace vectoriel $V$ sur $K$ (et en particulier pour
$V=L$). Or cela est évident si $V=K^n$ puisque alors
$W\otimes_KV=W^n$ pour tout espace vectoriel $W$ sur $K$ et que les produits
et limites projectives commutent. D'autre part, pour tout élément $z$ du
premier membre de (21.8.3.2), il existe un sous-espace $V'$ de $V$ de rang
fini tel que $z\in\Omega^1_{K/k_0}\otimes_KV'$ (identifié canoniquement à un
sous-espace de $\Omega^1_{K/k_0}\otimes_KV$)~; si $z\ne0$, son image dans
$\varprojlim_\alpha(\Omega^1_{K/k_\alpha}\otimes_KV')$ n'est donc pas nulle~;
comme le foncteur $\varprojlim$ est exact à gauche, l'image de $z$ dans le
second membre de (21.8.3.2) est donc aussi $\ne0$, ce qui achève de montrer
que $c)$ implique $b)$.

Il est trivial que $b)$ entraîne $b')$~; montrons que $b')$ entraîne $a)$.
Avec les notations de $b')$, on peut supposer que $L\ne K$. Il résulte alors
de (21.4.7) que l'on a $\operatorname{rg}\Upsilon_{L/K}\leqslant1$. Si
$\Upsilon_{L/K/k_\alpha}=0$, il n'y a rien à démontrer en vertu de
(21.6.1.1). Sinon, $\Upsilon_{L/K/k_\alpha}$ s'identifie canoniquement à
$\Upsilon_{L/K}$ et, si l'on pose $a=x^p$, $\Upsilon_{L/K}$ a une base
formée du seul élément $d_K(a)\otimes1$ (21.4.7)~;
$\Upsilon_{L/K/k_\alpha}$ a donc une base formée du seul élément
$d_{K/k_\alpha}(a)\otimes1$, et dire qu'un sous-corps $k_\alpha$ est
$k_0$-admissible pour l'extension $L$ de $K$ signifie que l'on a
$d_{K/k_\alpha}(a)\ne0$, ou encore (21.4.5) que
$a\notin k_\alpha(K^p)$. Or, pour tout $a\notin k_0(K^p)$, on a
$d_{K/k_0}(a)\ne0$ et $K(a^{1/p})\ne K$, donc on peut appliquer $b')$, qui
prouve l'existence d'un $\alpha$ tel que $a\notin k_\alpha(K^p)$~;
autrement dit $b')$ implique $a)$.

Reste à montrer que $a)$ implique $c)$. Soit $(x_\mu)_{\mu\in M}$ une
$p$-base de $K$ sur $k_0$~; alors, les $d_{K/k_0}(x_\mu)$ forment une base du
$K$-espace vectoriel $\Omega^1_{K/k_0}$ (21.4.5), et la condition $c)$
signifie que pour toute partie finie $J$ de l'ensemble d'indices $M$, il
existe un $\alpha\in I$ tel que les $d_{K/k_\alpha}(x_\mu)$ pour $\mu\in J$
soient linéairement indépendants dans $\Omega^1_{K/k_\alpha}$~; mais cela
signifie aussi (21.4.5) que les $x_\mu$ pour $\mu\in J$ forment une famille
$p$-libre sur $k_\alpha$, et l'existence d'un $\alpha$ ayant cette propriété
découle de (21.8.2) et de l'hypothèse $a)$.
\end{proof}
\oldpage[0\textsubscript{IV-1}]{174}

\begin{corollary}[21.8.4]
\label{0.21.8.4-fr}
Soient $K$ un corps de caractéristique $p>0$, $k_0$ un sous-corps de $K$,
$(k_\alpha)_{\alpha\in I}$ une famille filtrante décroissante de sous-corps
de $K$, contenant $k_0$, telle que
\[
  \bigcap_{\alpha\in I}k_\alpha(K^p)=k_0(K^p).
\]
Soient $L$ une extension de $K$ telle que $\Upsilon_{L/K/k_0}$ soit un
$L$-espace vectoriel de rang fini~; alors, pour tout corps $K'$ tel que
$K\subset K'\subset L$, il existe $\alpha\in I$ tel que $k_\alpha$ soit
$k_0$-admissible pour l'extension $L$ de $K'$.
\end{corollary}

\begin{proof}
Il suffit d'appliquer (21.8.3) et (21.6.6).
\end{proof}

\begin{corollary}[21.8.5]
\label{0.21.8.5-fr}
Soient $K$ un corps de caractéristique $p>0$, $k_0$ un sous-corps de $K$,
$(k_\alpha)_{\alpha\in I}$ une famille filtrante décroissante de sous-corps
de $K$ contenant $k_0$, telle que
\[
  \bigcap_{\alpha\in I}k_\alpha(K^p)=k_0(K^p).
\]
Alors, pour toute extension $L$ de $K$ telle que $\Upsilon_{L/K/k_0}$ soit
un $L$-espace vectoriel de rang fini, on a
\[
  \bigcap_{\alpha\in I}k_\alpha(L^p)=k_0(L^p).
\]
\end{corollary}

\begin{proof}
Supposons en effet que $M$ soit une extension de $L$ telle que
$\Upsilon_{M/L/k_0}$ soit un $M$-espace vectoriel de rang fini. La suite
exacte (21.6.1.1)
\[
  0\longrightarrow \Upsilon_{L/K/k_0}\otimes_L M
  \longrightarrow \Upsilon_{M/K/k_0}
  \longrightarrow \Upsilon_{M/L/k_0}
\]
et l'hypothèse montrent alors que $\Upsilon_{M/K/k_0}$ est aussi un
$M$-espace vectoriel de rang fini. Il existe donc un indice $\alpha$ tel que
$k_\alpha$ soit $k_0$-admissible pour l'extension $M$ de $K$ (21.8.3), donc
aussi pour l'extension $M$ de $L$ (21.6.6)~; cela ayant lieu pour toute
extension $M$ de $L$ telle que $\Upsilon_{M/L/k_0}$ soit de rang fini, le
corollaire résulte de l'équivalence de $a)$ et $b)$ dans (21.8.3).
\end{proof}

\begin{corollary}[21.8.6]
\label{0.21.8.6-fr}
Soient $K$ un corps, $p$ son exposant caractéristique, $k_0$ un sous-corps de
$K$. Si $L$ est une extension de $K$ telle que $\Upsilon_{L/K/k_0}$ soit un
$L$-espace vectoriel de rang fini, il existe un sous-corps $k$ de $K$,
contenant $k_0(K^p)$, tel que $[K:k]$ soit fini, et qui soit
$k_0$-admissible pour l'extension $L$ de $K$.
\end{corollary}

\begin{proof}
Il suffit en effet, en vertu de (21.8.3), de construire une famille filtrante
décroissante $(k_\alpha)_{\alpha\in I}$ de sous-corps de $K$, contenant
$K^p$ et $k_0$, pour lesquels $[K:k_\alpha]<+\infty$ et
$\bigcap_\alpha k_\alpha=k_0(K^p)$. Pour cela on considère une $p$-base
$(x_\lambda)_{\lambda\in J}$ de $K$ sur $k_0$ et, pour toute partie finie
$H$ de $J$, on considère le sous-corps $k_H$ de $K$ engendré par $k_0(K^p)$
et les $x_\lambda$ d'indice $\lambda\in J-H$~; il résulte de cette définition
que $(x_\lambda)_{\lambda\in H}$ est une $p$-base de $K$ sur $k_H$, et on en
conclut aussitôt que les $k_H$ vérifient les conditions voulues.
\end{proof}

\begin{remarks}[21.8.7]
\label{0.21.8.7-fr}
\begin{enumerate}
  \item[(i)] On a déjà vu (21.7.2) que si $L$ est une extension de
  \emph{type fini} de $K$, $\Upsilon_{L/K/k_0}$ est de rang fini pour tout
  sous-corps $k_0$ de $K$. Il en est de même si $L$ est une extension
  \emph{séparable} de $K$, car en vertu de (20.6.19), on a
  $\Upsilon_{L/K/k_0}=0$. Enfin, si $L$ est une \emph{extension de type fini
  d'une extension séparable $L_0$ de $K$}, le même raisonnement que dans
  (21.8.5) montre que $\Upsilon_{L/K/k_0}$ est encore de rang fini (et en fait
  est isomorphe à un sous-espace de $\Upsilon_{L/L_0/k_0}$).

  \oldpage[0\textsubscript{IV-1}]{175}
  \item[(ii)] Dans l'énoncé de (21.8.5), et par suite aussi dans celui de
  (21.8.3, $b)$), on ne peut omettre l'hypothèse que $\Upsilon_{L/K/k_0}$ est
  de rang fini sur $L$. Prenons pour $k_0$ le corps premier $\mathbf F_p$,
  pour $K$ un corps tel que $[K:K^p]$ soit infini dénombrable (par exemple le
  corps de fractions rationnelles $\mathbf F_p(X_1,\ldots,X_n,\ldots)$ à une
  infinité d'indéterminées)~; le procédé de construction de (21.8.6) montre
  aussitôt qu'il existe une suite infinie strictement décroissante $(K_n)$ de
  sous-corps de $K$ telle que $K_0=K$ et $\bigcap_nK_n=K^p$. Nous allons
  construire une suite croissante d'extensions finies $M_n$ de $K$, telle que
  si $M$ est la réunion des $M_n$, l'extension $L=M^{1/p}$ de $K$ mette
  (21.8.5) en défaut. Pour cela, soit $x$ un élément de $K-K^p$, posons
  $M_0=K^p$, et $M_n=K^p(a_1x,a_2x,\ldots,a_nx)$ pour $n\geqslant1$, où les
  $a_n$ sont construits par récurrence de façon que $a_n\in K_n$ et
  $a_{n+1}\notin M_n(x)$ pour tout $n$~: cela est possible, car $M_n(x)$ est
  de degré fini sur $K^p$, tandis qu'il n'en est pas de même de $K_n$. On en
  conclut aussitôt que $x\notin M=L^p$, mais comme
  $x=(a_nx)a_n^{-1}$, on a $x\in K_n(M)=K_n(L^p)$ pour tout $n$.
\end{enumerate}
\end{remarks}

\begin{proposition}[21.8.8]
\label{0.21.8.8-fr}
Soit $k$ un corps de caractéristique $p>0$, et soient
$A=k[[T_1,\ldots,T_r]]$ l'anneau des séries formelles à $r$ indéterminées sur
$k$, $K=k((T_1,\ldots,T_r))$ son corps des fractions. Alors il existe une
famille filtrante décroissante $(A_\alpha)$ de sous-anneaux noethériens de $A$
telle que $A$ soit un $A_\alpha$-module libre de type fini pour tout $\alpha$
et que, si $K_\alpha$ est le corps des fractions de $A_\alpha$, on ait
$\bigcap_\alpha K_\alpha=K^p$.
\end{proposition}

\begin{proof}
On peut écrire $K^p=k^p((T_1^p,\ldots,T_r^p))$~; on a vu dans la démonstration
de (21.8.6) qu'il existe une famille décroissante $(k_\alpha)$ de sous-corps
de $k$ telle que $[k:k_\alpha]$ soit fini pour tout $\alpha$ et
$\bigcap_\alpha k_\alpha=k^p$~; il est clair que si l'on pose
$A_\alpha=k_\alpha[[T_1^p,\ldots,T_r^p]]$, $A$ est un $A_\alpha$-module
libre de type fini~; tout revient donc à prouver la relation
\begin{equation}
\tag{21.8.8.1}
\label{0.21.8.8.1-fr}
  \bigcap_\alpha k_\alpha((T_1^p,\ldots,T_r^p))
  =k^p((T_1^p,\ldots,T_r^p)).
\end{equation}
Comme $\bigcap_\alpha k_\alpha=k^p$, cela va résulter des deux lemmes suivants.
\end{proof}

\begin{lemma}[21.8.8.2]
\label{0.21.8.8.2-fr}
Si $k'$ est une extension d'un corps $k$, on a
\[
  k'[[T_1,\ldots,T_r]]\cap k((T_1,\ldots,T_r))=k[[T_1,\ldots,T_r]].
\]
\end{lemma}

\begin{proof}
En effet, posons $C=k[[T_1,\ldots,T_r]]$, $D=k'[[T_1,\ldots,T_r]]$~; comme
$k((T_1,\ldots,T_r))$ est le corps des fractions de $C$, il suffira de
prouver que $D$ est un $C$-module \emph{fidèlement plat} (Bourbaki,
\emph{Alg. comm.}, chap.~I, \S~3, n\textsuperscript{o}~5, prop.~10). Or, $C$
et $D$ sont des anneaux locaux noethériens, et si $\mathfrak m$ est l'idéal
maximal de $C$, on a
$D/\mathfrak m^jD=(C/\mathfrak m^j)\otimes_k k'$, donc
$D/\mathfrak m^jD$ est un $(C/\mathfrak m^j)$-module plat~; il suffit donc
d'appliquer ($\mathbf 0_{\mathrm{III}}$, 10.2.1) et
($\mathbf 0_{\mathrm I}$, 6.6.2).
\end{proof}

\begin{lemma}[21.8.8.3]
\label{0.21.8.8.3-fr}
Soient $k$ un corps, $(k_\alpha)$ une famille filtrante décroissante de
sous-corps de $k$, et posons $k_0=\bigcap_\alpha k_\alpha$. On suppose qu'il
existe une puissance $q$ de l'exposant caractéristique de $k$ telle que
$k^q\subset k_0$. Alors on a
\begin{equation}
\tag{21.8.8.4}
\label{0.21.8.8.4-fr}
  \bigcap_\alpha k_\alpha((T_1,\ldots,T_r))=k_0((T_1,\ldots,T_r)).
\end{equation}
\end{lemma}

\begin{proof}
\oldpage[0\textsubscript{IV-1}]{176}
Il suffit de prouver qu'un élément $f\ne0$ du premier membre de (21.8.8.4)
appartient au second membre. Soient $\gamma$ un indice,
$g\in k_\gamma[[T_1,\ldots,T_r]]$ tel que
$gf\in k_\gamma[[T_1,\ldots,T_r]]$~; quitte à remplacer $g$ par $g^q$, on
peut supposer que $g\in k_0[[T_1,\ldots,T_r]]$. Alors, pour tout
$\alpha\geqslant\gamma$, on a
\[
  gf\in k_\gamma[[T_1,\ldots,T_r]]\cap k_\alpha((T_1,\ldots,T_r))
  =k_\alpha[[T_1,\ldots,T_r]]
\]
en vertu du lemme (21.8.8.2). Mais il est clair que l'intersection des
anneaux $k_\alpha[[T_1,\ldots,T_r]]$ n'est autre que
$k_0[[T_1,\ldots,T_r]]$, et l'on a donc bien
$f\in k_0((T_1,\ldots,T_r))$.
\end{proof}

\subsection{Modules de différentielles complétés dans les anneaux de séries formelles}
\label{subsection:0.21.9-fr}

Dans ce numéro, les corps ne sont plus nécessairement supposés être de
caractéristique $>0$.

\begin{lemma}[21.9.1]
\label{0.21.9.1-fr}
Soient $k$ un corps, $A$ un anneau local noethérien complet qui est une
$k$-algèbre, $K$ le corps résiduel de $A$. Si $K$ est une extension de type
fini de $k$, le $A$-module $\widehat\Omega^1_{A/k}$ est de type fini.
\end{lemma}

\begin{proof}
En vertu de (20.7.15), il suffit de prouver que $\Omega^1_{K/k}$ est un
$K$-espace vectoriel de rang fini. Or, par hypothèse, $K$ est le corps des
fractions d'une $k$-algèbre de type fini $B$~; comme $\Omega^1_{B/k}$ est un
$B$-module de type fini en vertu de (20.4.7), la conclusion résulte de
(20.5.9).
\end{proof}

\begin{proposition}[21.9.2]
\label{0.21.9.2-fr}
Soient $k_0$ un corps, $A$ un anneau local noethérien complet qui est une
$k_0$-algèbre formellement lisse, $\mathfrak m$ son idéal maximal,
$K=A/\mathfrak m$ son corps résiduel (de sorte qu'en tant qu'anneau $A$ est
isomorphe à un anneau de séries formelles $K[[T_1,\ldots,T_r]]$ (19.6.5)).
Si $K$ est une extension de type fini de $k_0$, alors
$\widehat\Omega^1_{A/k_0}$ est un $A$-module libre de rang égal à
$\dim(A)+\deg.\operatorname{tr}_{k_0}K$. En outre, pour tout sous-corps $k$
de $k_0$ tel que $\Omega^1_{k_0/k}$ soit de rang fini,
$\widehat\Omega^1_{A/k}$ est un $A$-module libre de rang égal à
\[
  \dim(A)+\deg.\operatorname{tr}_{k_0}K
  +\operatorname{rg}_{k_0}(\Omega^1_{k_0/k}).
\]
\end{proposition}

\begin{proof}
Il est clair que, si $P$ est le sous-corps premier de $k_0$, $K$ est une
$P$-algèbre formellement lisse (19.6.1). Notons d'autre part que l'on a
$\Upsilon^K_{A/k_0/P}=0$~: en effet, cela revient à voir que l'homomorphisme
\[
  u_{A/k_0/P}\otimes1_K:
  \Omega^1_{k_0/P}\otimes_{k_0}K
  \longrightarrow \Omega^1_{A/P}\otimes_A K
\]
est injectif (20.6.14.5), $u$ désignant l'homomorphisme structural
$k_0\to A$. Mais comme $A$ est une $k_0$-algèbre formellement lisse par
hypothèse, il résulte de (20.7.2) que $u_{A/k_0/P}$ est formellement
inversible à gauche, donc (19.1.7) $u_{A/k_0/P}\otimes1_K$ est inversible à
gauche et \emph{a fortiori} injectif, ce qui prouve notre assertion.
Appliquant la suite exacte (20.6.22.1), où $\Lambda,A,B,C$ sont remplacés par
$P,k_0,A,K$, il vient la suite exacte
\[
  0\longrightarrow \Upsilon^K_{A/k_0}
  \longrightarrow \mathfrak m/\mathfrak m^2
  \longrightarrow \Omega^1_{A/k_0}\otimes_A K
  \longrightarrow \Omega^1_{K/k_0}\longrightarrow0,
\]
de sorte que l'on a
\[
\begin{split}
  \operatorname{rg}_K(\Omega^1_{A/k_0}\otimes_AK)
  &={\operatorname{rg}}_K(\mathfrak m/\mathfrak m^2)
    +\bigl(\operatorname{rg}_K(\Omega^1_{K/k_0})
    -\operatorname{rg}_K(\Upsilon_{K/k_0})\bigr)\\
  &={\operatorname{rg}}_K(\mathfrak m/\mathfrak m^2)
    +\deg.\operatorname{tr}_{k_0}K
\end{split}
\]
en vertu de (21.7.1). D'autre part, $\widehat\Omega^1_{A/k_0}$ est un
$A$-module formellement projectif (20.4.9), et comme sa topologie est la
topologie $\mathfrak m$-préadique (20.4.5),
$\widehat\Omega^1_{A/k_0}/\mathfrak m^{j+1}\widehat\Omega^1_{A/k_0}$ est un
$(A/\mathfrak m^{j+1})$-module projectif pour tout $j$ (19.2.4), donc libre
($\mathbf 0_{\mathrm{III}}$, 10.1.2) et de rang

\oldpage[0\textsubscript{IV-1}]{177}
$n=\operatorname{rg}_K(\Omega^1_{A/k_0}\otimes_AK)$. En vertu de (21.9.1),
$\widehat\Omega^1_{A/k_0}$ est un $A$-module de type fini, et en outre ce
$A$-module est plat en vertu de ($\mathbf 0_{\mathrm{III}}$, 10.2.1)~; il est
donc libre ($\mathbf 0_{\mathrm{III}}$, 10.1.3) de rang $n$, d'où la première
assertion.

Comme par hypothèse $\Omega^1_{k_0/k}$ est de rang fini sur $k_0$, le produit
tensoriel complété
$\Omega^1_{k_0/k}\widehat\otimes_{k_0}A$ s'identifie à
$\Omega^1_{k_0/k}\otimes_{k_0}A$~; comme $A$ est une $k_0$-algèbre
formellement lisse, il résulte de (20.7.18) que l'homomorphisme
\[
  \widehat u_{A/k_0/k}:
  \Omega^1_{k_0/k}\otimes_{k_0}A
  \longrightarrow\widehat\Omega^1_{A/k}
\]
admet un inverse à gauche qui est un $A$-homomorphisme continu, ce qui
implique en particulier que son image est fermée dans
$\widehat\Omega^1_{A/k}$~; la suite exacte
\[
  0\longrightarrow \Omega^1_{k_0/k}\otimes_{k_0}A
  \xrightarrow{\widehat u_{A/k_0/k}}\widehat\Omega^1_{A/k}
  \longrightarrow\widehat\Omega^1_{A/k_0}\longrightarrow0
\]
est donc exacte et scindée en vertu de (20.7.17) et de ce qui précède~; ce
qui achève de prouver la proposition.
\end{proof}

\begin{corollary}[21.9.3]
\label{0.21.9.3-fr}
Soient $k$ un corps, $A$ l'anneau de séries formelles
$k[[T_1,\ldots,T_r]]$, muni de sa structure usuelle de $k$-algèbre. Alors
$\widehat\Omega^1_{A/k}$ est un $A$-module libre de rang $r=\dim(A)$.
\end{corollary}

\begin{proof}
Il suffit de noter que $A$ est une $k$-algèbre formellement lisse (19.3.4),
et d'appliquer (21.9.2) avec $k_0=k=K$.
\end{proof}

\begin{lemma}[21.9.4]
\label{0.21.9.4-fr}
Soient $k$ un corps de caractéristique $p>0$, $A$ une $k$-algèbre qui est un
anneau local noethérien complet, $B$ une sous-$k$-algèbre de $A$ isomorphe à
une $k$-algèbre topologique de la forme $k[[T_1,\ldots,T_r]]$ et telle que
$A$ soit une $B$-algèbre de type fini. Si $B_1$ est la sous-algèbre
$k[[T_1^p,\ldots,T_r^p]]$ de $B$, $\widehat\Omega^1_{A/k}$ s'identifie
canoniquement à $\Omega^1_{A/B_1}$.
\end{lemma}

\begin{proof}
Toute dérivation continue de $B_1$ dans un $A$-module topologique $P$, qui
est restriction d'une $k$-dérivation continue de $A$ dans $P$, est \emph{nulle},
car elle l'est dans le sous-anneau $k[T_1^p,\ldots,T_r^p]$ des polynômes, et
ce dernier est dense dans $B_1$. La suite exacte (20.3.7.1) montre donc que
l'homomorphisme canonique
\[
  \operatorname{D\acute{e}r.\,cont}_{B_1}(A,P)
  \longrightarrow \operatorname{D\acute{e}r.\,cont}_k(A,P)
\]
est bijectif~; la conclusion résulte de (20.7.14.4), compte tenu de ce que
$\Omega^1_{A/B_1}$ est un $A$-module de type fini (20.4.7), donc séparé et
complet puisque $A$ est un anneau local noethérien complet.
\end{proof}

\begin{proposition}[21.9.5]
\label{0.21.9.5-fr}
Soient $A$ un anneau local noethérien complet intègre, $A_0$ un sous-anneau
de $A$ isomorphe à un anneau de séries formelles $k[[T_1,\ldots,T_r]]$ sur
un corps $k$, tel que $A$ soit une $A_0$-algèbre finie et que le corps des
fractions $E$ de $A$ soit une extension séparable du corps des fractions
$L_0$ de $A_0$. Alors on a
\begin{equation}
\tag{21.9.5.1}
\label{0.21.9.5.1-fr}
  \operatorname{rg}_A\widehat\Omega^1_{A/k}
  =\operatorname{rg}_E(\widehat\Omega^1_{A/k}\otimes_AE)=\dim(A).
\end{equation}
\end{proposition}

\begin{proof}
Notons que si $\mathfrak m_0$ est l'idéal maximal de $A_0$, la topologie de
$A$ est la topologie $\mathfrak m_0$-adique puisque $A$ est une
$A_0$-algèbre finie (Bourbaki, \emph{Alg. comm.}, chap.~IV, \S~2,
n\textsuperscript{o}~5, cor.~3 de la prop.~9) et induit sur $A_0$ la topologie
$\mathfrak m_0$-adique (Bourbaki, \emph{Alg. comm.}, chap.~III, \S~3,
n\textsuperscript{o}~4, th.~3). On sait (21.9.1) que
$\widehat\Omega^1_{A/k}$ est un $A$-module de type fini et,

\oldpage[0\textsubscript{IV-1}]{178}
en vertu de (21.9.3) $\widehat\Omega^1_{A_0/k}$ est un $A_0$-module libre de
rang $r=\dim(A_0)=\dim(A)$ (16.1.5)~; le produit tensoriel
$\widehat\Omega^1_{A_0/k}\otimes_{A_0}A$ est donc un $A$-module libre de rang
$r$ identique à son séparé complété~; enfin $\widehat\Omega^1_{A/A_0}$ est
un $A$-module de type fini (20.4.7), donc identique à son séparé complété
($\mathbf 0_{\mathrm I}$, 7.3.6). Cela étant, dans la suite d'homomorphismes
\begin{equation}
\tag{21.9.5.2}
\label{0.21.9.5.2-fr}
  \widehat\Omega^1_{A_0/k}\otimes_{A_0}A
  \xrightarrow{u}\widehat\Omega^1_{A/k}
  \xrightarrow{v}\widehat\Omega^1_{A/A_0}\longrightarrow0
\end{equation}
on sait que $v$ est surjectif et que $\operatorname{Im}(u)$ est dense dans
$\operatorname{Ker}(v)$ (20.7.17)~; mais tout sous-$A_0$-module de
$\widehat\Omega^1_{A/k}$ est fermé pour la topologie $\mathfrak m_0$-adique
($\mathbf 0_{\mathrm I}$, 7.3.5), donc la suite (21.9.5.2) est exacte. Notons
d'autre part que puisque $A$ est entier sur $A_0$, $A_0$ intégralement clos
et $E$ séparable sur $L_0$, $\widehat\Omega^1_{A/A_0}$ est un $A$-module de
torsion (20.4.13, (iv))~; tensorisant la suite exacte (21.9.5.2) par $E$, il
vient une suite exacte
\[
  \widehat\Omega^1_{A_0/k}\otimes_{A_0}E
  \longrightarrow \widehat\Omega^1_{A/k}\otimes_AE\longrightarrow0
\]
d'où $\operatorname{rg}_E(\widehat\Omega^1_{A/k}\otimes_AE)\leqslant r$. Il
reste donc à montrer qu'il existe, dans $\widehat\Omega^1_{A/k}\otimes_AE$,
$r$ éléments linéairement indépendants sur $E$. Or, soient
$D_i=\partial/\partial T_i$ les dérivations canoniques de $A_0$ dans lui-même
$(1\leqslant i\leqslant r)$~; elles se prolongent de façon unique en des
dérivations (encore notées $D_i$) de $E$ dans lui-même, puisque $E$ est
extension séparable finie de $L_0$~; par restriction à $A$, ces dérivations
donnent des dérivations de $A$ dans $E$, et il est immédiat que ces
dérivations prennent leurs valeurs dans un même sous-$A$-module de type fini
de $E$~; comme elles sont continues dans $A_0$ et que la topologie de $A$ est
la topologie $\mathfrak m_0$-adique, on a ainsi formé $r$ dérivations
continues de $A$ dans un $A$-module topologique, qui sont évidemment
linéairement indépendantes puisque $D_i(T_j)=\delta_{ij}$~; ceci achève de
prouver la proposition.
\end{proof}

La proposition suivante est un analogue «~formel~» de (21.7.3)~:

\begin{proposition}[21.9.6]
\label{0.21.9.6-fr}
Soient $k_0\subset k\subset K_0$ trois corps de caractéristique $p>0$, tels
que $[K_0:k]<+\infty$~; on pose
\[
  L_0=K_0((T_1,\ldots,T_r)),\qquad
  M_0=K_0((T_1^p,\ldots,T_r^p)),
\]
\[
  L=k((T_1,\ldots,T_r)),\qquad
  M=k((T_1^p,\ldots,T_r^p)),\qquad
  N=k((T_1^{p^2},\ldots,T_r^{p^2})).
\]
Soit $E$ une extension de type fini de $L_0$.
\begin{enumerate}
  \item[(i)] Si $\Upsilon_{K_0/k_0}$ est de rang fini, on a
  \begin{equation}
  \tag{21.9.6.1}
  \label{0.21.9.6.1-fr}
    \operatorname{rg}_E\Omega^1_{E/M}
    -\operatorname{rg}_{K_0}\Omega^1_{K_0/k}
    =(r+\deg.\operatorname{tr}_{L_0}E)
      +d(E/K_0,k_0)+d(E/M_0,N/k_0)
  \end{equation}
  (notations de (21.6.1)), et en particulier on a
  \begin{equation}
  \tag{21.9.6.2}
  \label{0.21.9.6.2-fr}
    \operatorname{rg}_E\Omega^1_{E/M}
    -\operatorname{rg}_{K_0}\Omega^1_{K_0/k}
    \geqslant r+\deg.\operatorname{tr}_{L_0}E+d(E/K_0,k_0).
  \end{equation}
  En outre, si les deux membres de (21.9.6.2) sont égaux, ils le restent
  lorsqu'on remplace $k$ par un corps $k'$ tel que $k_0\subset k'\subset k$.

  \item[(ii)] On suppose en outre que $[k_0:k_0^p]<+\infty$. Soit
  $(k_\alpha)$ une famille filtrante décroissante de sous-corps de $K_0$
  contenant $k_0$, tels que $[K_0:k_\alpha]<+\infty$ pour tout $\alpha$ et
  que
  $\bigcap_\alpha k_\alpha(K_0^p)=k_0(K_0^p)$~; alors il existe un $\alpha$
  tel que, pour $k=k_\alpha$, les deux membres de (21.9.6.2) soient égaux.
\end{enumerate}
\end{proposition}

\begin{proof}
\begin{enumerate}
  \oldpage[0\textsubscript{IV-1}]{179}
  \item[(i)] On sait (21.1.5, (ii)) que
  $\operatorname{rg}_E\Omega^1_{E/M}$ ne change pas quand on remplace $M$ par
  $M(L_0^p)$~; comme
  $L_0^p=K_0^p((T_1^p,\ldots,T_r^p))$ et que
  $[k(K_0^p):k]<+\infty$, on a
  \[
    M(L_0^p)=k(K_0^p)((T_1^p,\ldots,T_r^p))~;
  \]
  de même $\Omega^1_{K_0/k}$ ne change pas quand on y remplace $k$ par
  $k(K_0^p)$, si bien que l'on peut \emph{supposer que $K_0^p\subset k$}, ce
  que nous ferons dans toute la suite. Nous introduirons aussi les corps
  \[
    k'_0=k_0(K_0^p),\qquad
    N'=k'_0((T_1^{p^2},\ldots,T_r^{p^2}))
  \]
  de sorte que l'on a le diagramme de corps
  \[
    \xymatrix{
      K_0\ar[r] & M_0\ar[r] & L_0\ar[r] & E\\
      k\ar[u]\ar[r] & N\ar[u]\ar[r] & M\ar[u]\ar[r] & L\ar[u]\\
      k_0\ar[u]\ar[r] & k'_0\ar[u]\ar[r] & N'\ar[u] &
    }
  \]

  Nous nous proposons d'évaluer la différence $\delta$ du premier et du
  second membre de (21.9.6.2).

  Il résulte de (21.7.3) que l'on a
  \[
    \operatorname{rg}_E\Omega^1_{E/M}
    -\operatorname{rg}_{M_0}\Omega^1_{M_0/M}
    =\deg.\operatorname{tr}_{M_0}E+d(E/M_0,M)
  \]
  et comme $L_0$ est une extension finie de $M_0$, on a
  $\deg.\operatorname{tr}_{L_0}E=\deg.\operatorname{tr}_{M_0}E$. D'autre
  part, une $p$-base de $K_0$ sur $k=k(K_0^p)$ est aussi une $p$-base de
  $M_0$ sur $M$, donc (21.4.5), on a
  \[
    \operatorname{rg}_{M_0}\Omega^1_{M_0/M}
    =\operatorname{rg}_{K_0}\Omega^1_{K_0/k}
  \]
  si bien que l'on a
  \begin{equation}
  \tag{21.9.6.3}
  \label{0.21.9.6.3-fr}
    \delta=d(E/M_0,M)-d(E/K_0,k_0)-r.
  \end{equation}

  Notons maintenant le lemme classique~:

  \begin{lemma}[21.9.6.4]
  \label{0.21.9.6.4-fr}
  Pour tout corps $k$, le corps de séries formelles
  $K=k((T_1,\ldots,T_r))$ est une extension séparable de $k$.
  \end{lemma}

  \begin{proof}
  Rappelons rapidement la démonstration de ce lemme pour être complet. Il
  suffit (Bourbaki, \emph{Alg.}, chap.~VIII, \S~7,
  n\textsuperscript{o}~3, démonstration du th.~1) de prouver que pour toute
  extension finie $k'$ de $k$, $K\otimes_k k'$ est sans élément nilpotent~;
  mais si l'on pose $A=k[[T_1,\ldots,T_r]]$ et $S=A-\{0\}$,
  $K\otimes_k k'$ est égal à $S^{-1}(A\otimes_k k')$, et
  $A\otimes_k k'$ s'identifie canoniquement à l'anneau intègre
  $A'=k'[[T_1,\ldots,T_r]]$ et $A$ à un sous-anneau de $A'$, d'où la
  conclusion.
  \end{proof}

  Utilisant ce lemme et (21.6.2, (ii)), on a $d(L_0/K_0,k_0)=0$, et la
  formule (21.6.4.4) donne donc
  \[
    d(E/K_0,k_0)=d(E/L_0,k_0).
  \]

  \oldpage[0\textsubscript{IV-1}]{180}
  Les formules (21.6.4.4) et (21.6.4.2) donnent d'autre part
  \[
    d(E/M_0,M)=d(E/L_0,M)+d(L_0/M_0,M),
  \]
  \[
    d(L_0/M_0,M)=d(L_0/M_0,M/N)+d(L_0/M_0,N),
  \]
  d'où
  \[
    \delta=d(E/L_0,M)+d(L_0/M_0,N)-d(E/L_0,k_0)
    +d(L_0/M_0,M/N)-r.
  \]

  Nous allons montrer que l'on a la relation
  \begin{equation}
  \tag{21.9.6.5}
  \label{0.21.9.6.5-fr}
    d(L_0/M_0,M/N)=r.
  \end{equation}
  Pour cela, remarquons que, par définition, on a
  \[
    d(L_0/M_0,M/N)
    =\operatorname{rg}_{L_0}\Upsilon_{L_0/M/N}
    -\operatorname{rg}_{M_0}\Upsilon_{M_0/M/N}.
  \]
  Compte tenu des deux suites exactes
  \[
  \begin{aligned}
    0&\longrightarrow\Upsilon_{L_0/M/N}
      \longrightarrow\Omega^1_{M/N}\otimes_ML_0
      \longrightarrow\Omega^1_{L_0/N}
      \longrightarrow\Omega^1_{L_0/M}\longrightarrow0,\\
    0&\longrightarrow\Upsilon_{M_0/M/N}
      \longrightarrow\Omega^1_{M/N}\otimes_MM_0
      \longrightarrow\Omega^1_{M_0/N}
      \longrightarrow\Omega^1_{M_0/M}\longrightarrow0,
  \end{aligned}
  \]
  il vient
  \[
  \begin{split}
    d(L_0/M_0,M/N)={}&\operatorname{rg}_{L_0}\Omega^1_{L_0/M}
      -\operatorname{rg}_{L_0}\Omega^1_{L_0/N}\\
      &-\operatorname{rg}_{M_0}\Omega^1_{M_0/M}
      +\operatorname{rg}_{M_0}\Omega^1_{M_0/N}.
  \end{split}
  \]
  Or, on a $N(L_0^p)\supset M$, donc
  $\Omega^1_{L_0/N}=\Omega^1_{L_0/M}$ (21.1.5, (ii)). D'autre part, les
  $T_i^p$ $(1\leqslant i\leqslant r)$ forment une $p$-base de $M$ sur $N$,
  et une $p$-base de $K_0$ sur $k$ est aussi une $p$-base de $M_0$ sur $M$,
  donc, comme $M_0^p\subset N$,
  \[
    \operatorname{rg}_{M_0}\Omega^1_{M_0/N}
    =r+\operatorname{rg}_{K_0}\Omega^1_{K_0/k}
    =r+\operatorname{rg}_{M_0}\Omega^1_{M_0/M}
  \]
  en vertu de (21.4.5) et (21.1.10), ce qui prouve (21.9.6.5).

  On a encore, par (21.6.4.2),
  \[
    d(E/L_0,M)=d(E/L_0,k_0)+d(E/L_0,M/k_0)
  \]
  et
  \[
    d(L_0/M_0,N)=d(L_0/M_0,N/k)+d(L_0/M_0,k).
  \]
  Mais comme $L_0$ est séparable sur $K_0$ (21.9.6.4), on a
  $d(L_0/K_0,k)=0$, et par suite aussi $d(L_0/M_0,k)=0$ ((21.6.2) et
  (21.6.4)). Il vient donc
  \[
    \delta=d(E/L_0,M/k_0)+d(L_0/M_0,N/k).
  \]
  Appliquons encore deux fois (21.6.4.2) au premier terme du second membre de
  cette formule~; on obtient
  \[
    d(E/L_0,M/k_0)=d(E/L_0,M/N)+d(E/L_0,N/k)+d(E/L_0,k/k_0).
  \]
  Comme $M\subset N(L_0^p)$, on a $d(E/L_0,M/N)=0$ (21.6.2), donc (par
  (21.6.4.4))
  \[
    \delta=d(E/L_0,k/k_0)+d(E/M_0,N/k).
  \]
  Mais $M_0$ et $L_0$ sont des extensions séparables de $K_0$ (21.9.6.4),
  donc
  \[
    d(L_0/K_0,k/k_0)=d(M_0/K_0,k/k_0)=0
  \]
  (21.6.2), et appliquant encore deux fois (21.6.4.4),

  \oldpage[0\textsubscript{IV-1}]{181}
  on a
  \[
    d(E/L_0,k/k_0)=d(E/K_0,k/k_0)=d(E/M_0,k/k_0),
  \]
  d'où, par une dernière application de (21.6.4.2),
  \[
    \delta=d(E/M_0,N/k_0),
  \]
  ce qui prouve (21.9.6.1). En outre, lorsque $k$ est remplacé par une
  sous-extension $k'$ de $k_0$, $N$ est remplacé par une sous-extension $N'$,
  donc $d(E/M_0,N'/k_0)\leqslant d(E/M_0,N/k_0)$ (21.6.4.2), ce qui démontre
  la dernière assertion de (i).

  \item[(ii)] Pour tout $\alpha$, posons
  $N_\alpha=k_\alpha((T_1^{p^2},\ldots,T_r^{p^2}))$. En vertu de (21.9.6.1)
  et de (21.8.3), il s'agit de montrer (compte tenu de ce que
  $K_0^p\subset k_\alpha$ et de ce que $E$ est une extension de type fini de
  $M_0$) que l'on a
  \[
    \bigcap_\alpha N_\alpha=k_0(M_0^p).
  \]
  Or, on a
  \begin{equation}
  \tag{21.9.6.6}
  \label{0.21.9.6.6-fr}
    \bigcap_\alpha N_\alpha=N'
    =(k_0(K_0^p))((T_1^{p^2},\ldots,T_r^{p^2}))
  \end{equation}
  en vertu de (21.8.8.3) et de la relation
  $\bigcap_\alpha k_\alpha=k_0(K_0^p)$ (21.8.6). D'autre part, on a
  $k_0(M_0^p)=k_0(K_0^p((T_1^{p^2},\ldots,T_r^{p^2})))$. Pour achever la
  démonstration de (21.9.6, (ii)), il reste donc à prouver le

  \begin{lemma}[21.9.6.7]
  \label{0.21.9.6.7-fr}
  Soient $k_0$ un corps de caractéristique $p>0$ tel que
  $[k_0:k_0^p]<+\infty$, $K_0$ une extension de $k_0$. Alors on a
  \begin{equation}
  \tag{21.9.6.8}
  \label{0.21.9.6.8-fr}
    (k_0(K_0^p))((T_1,\ldots,T_r))
    =k_0\bigl(K_0^p((T_1,\ldots,T_r))\bigr).
  \end{equation}
  \end{lemma}

  \begin{proof}
  Comme $k_0^p\subset K_0^p$, on a
  \[
    k_0^p\bigl(K_0^p((T_1,\ldots,T_r))\bigr)
    =K_0^p((T_1,\ldots,T_r)).
  \]
  Si $(c_i)_{1\leqslant i\leqslant n}$ est une base de $k_0$ sur $k_0^p$,
  il est immédiat que $(c_i)$ est aussi un système de générateurs de chacun
  des deux membres de (21.9.6.8) sur $K_0^p((T_1,\ldots,T_r))$.
  \end{proof}
\end{enumerate}
\end{proof}

\begin{remark}[21.9.7]
\label{0.21.9.7-fr}
L'assertion de (21.9.6, (ii)) \emph{ne s'étend pas} au cas où
$[k_0:k_0^p]$ est infini. Prenons en effet pour $k_0$ le corps
$\mathbf F_p(X_n)_{n\geqslant0}$ des fractions rationnelles à une infinité
d'indéterminées, de sorte que les $X_n$ forment une $p$-base de $k_0$ sur
$\mathbf F_p$, et par suite $[k_0:k_0^p]=+\infty$. Dans l'énoncé de (21.9.6),
prenons $K_0=k_0$ (donc nécessairement $k=k_0$), $L_0=k_0((T))$, pour $E$
l'extension $L_0(y)$, où $y$ est racine du polynôme
\[
  Y^p-\sum_{n=0}^{\infty}X_nT^{np}
\]
de $L_0[Y]$. Alors la différence $\delta$ des deux membres de (21.9.6.2) est
\emph{non nulle}. En effet, $E$ est \emph{séparable} sur $k_0$, sans quoi
$y^p$ serait élément de $k_0(L_0^p)$, et l'on voit aussitôt que cela n'est
pas possible (dû au fait qu'il y a une infinité de $X_n$ algébriquement
indépendants)~; on a donc $d(E/K_0,k_0)=0$, la formule (21.9.6.3) se réduit à
$\delta=d(E/M_0,M_0)-1$, et il est clair que
$d(E/M_0,M_0)=\operatorname{rg}_E\Upsilon_{E/M_0}$. Tout revient à vérifier
que ce rang est $2$. Or on a $E^p\subset M_0$, une $p$-base de $E$ sur $L_0$
est formée du seul élément $y$, et $L_0$ a évidemment une $p$-base sur $M_0$
formée du seul élément $T$~; donc, puisque $E$ est extension algébrique de
$M_0$, notre assertion résulte de (21.7.1) et (21.4.5).
\end{remark}

\begin{proposition}[21.9.8]
\label{0.21.9.8-fr}
Soient $k_0$ un corps (de caractéristique quelconque), $K_0$ une extension
séparable de $k_0$, $A$ un anneau local noethérien complet, qui est une
$k_0$-algèbre et dont le corps résiduel

\oldpage[0\textsubscript{IV-1}]{182}
est une extension finie de $K_0$. Pour tout idéal premier $\mathfrak p$ de
$A$, soit $k(\mathfrak p)$ le corps résiduel de $A$ en $\mathfrak p$. Alors,
pour tout corps $k$ tel que $k_0\subset k\subset K_0$ et tel que
$[K_0:k]<+\infty$, on a
\begin{equation}
\tag{21.9.8.1}
\label{0.21.9.8.1-fr}
\begin{split}
  &\operatorname{rg}_{k(\mathfrak p)}
    \bigl(\widehat\Omega^1_{(A/\mathfrak p)/k}
      \otimes_{A/\mathfrak p}k(\mathfrak p)\bigr)
    -\operatorname{rg}_{K_0}\Omega^1_{K_0/k}\\
  &\hspace{35mm}\geqslant
    \dim(A/\mathfrak p)
    +\operatorname{rg}_{k(\mathfrak p)}\Upsilon_{k(\mathfrak p)/k_0}.
\end{split}
\end{equation}
Supposons en outre que $[k_0:k_0^p]<+\infty$ (où $p$ est l'exposant
caractéristique de $k_0$). Alors on peut trouver un corps $k$ tel que
$k_0(K_0^p)\subset k\subset K_0$ et $[K_0:k]<+\infty$, pour lequel les deux
membres de (21.9.8.1) sont égaux.
\end{proposition}

\begin{proof}
Comme $A/\mathfrak p$ a même corps résiduel $K$ que $A$ et est complet, on
peut se borner au cas où $A$ est intègre et $\mathfrak p=0$~; on posera
alors $E=k(\mathfrak p)$, corps des fractions de $A$. Comme $K_0$ est
formellement lisse sur $k_0$ (19.6.1), il existe un $k_0$-monomorphisme
$K_0\to A$ faisant de $A$ une $K_0$-algèbre~; on sait en outre qu'il existe
une sous-$K_0$-algèbre $A_0$ de $A$, $K_0$-isomorphe à une algèbre de séries
formelles $K_0[[T_1,\ldots,T_r]]$ et telle que $A$ soit une $A_0$-algèbre
\emph{finie} (19.8.9), ce qui entraîne que $E$ est une extension finie du
corps des fractions $L_0=K_0((T_1,\ldots,T_r))$ de $A_0$. On a alors
$\Upsilon_{K_0/k_0}=0$ puisque $K_0$ est séparable sur $k_0$, et
$\deg.\operatorname{tr}_{L_0}E=0$~; en outre, comme $A$ est une
$A_0$-algèbre finie on a $\dim(A)=\dim(A_0)$ (16.1.5).

Si $p=1$, on a
$\operatorname{rg}_E(\widehat\Omega^1_{A/k}\otimes_AE)=r$ par (21.9.5) et
(17.1.4, (iii)), et $\Omega^1_{K_0/k}=0$ puisque $K_0$ est extension finie
séparable de $k$ (21.7.4)~; d'autre part $\Upsilon_{E/k_0}=0$, $E$ étant
extension séparable de $k_0$~; les deux membres de (21.9.8.1) sont donc égaux
dans ce cas.

Supposons maintenant que $p>1$~; si $B=k[[T_1,\ldots,T_r]]$, $A$ est une
$B$-algèbre finie, donc, en posant $B_1=k[[T_1^p,\ldots,T_r^p]]$,
$\widehat\Omega^1_{A/k}$ s'identifie à $\Omega^1_{A/B_1}$ (21.9.4), et en
désignant par $M$ le corps des fractions $k((T_1^p,\ldots,T_r^p))$, il
résulte de (20.5.9) que $\widehat\Omega^1_{A/k}\otimes_AE$ s'identifie à
$\Omega^1_{E/M}$~; l'inégalité (21.9.8.1) n'est autre alors que (21.9.6.2),
et la dernière assertion du corollaire résulte de (21.9.6, (ii)).
\end{proof}
